\section*{Lecture 2: Information and Decision Modeling}

\addcontentsline{toc}{section}{Lecture 2: Information and Decision Modeling}

\clearpage

\SlideHeader{Lecture 2: Information and Decision Modeling}{001}{表紙}

\SlideImage{1}{../Materials/2026/3DM_02_Information-and-Decision-Modeling_SS26.pdf}

\vspace{1.5mm}

\begin{multicols}{2}

\footnotesize

\raggedcolumns

\NoteSection{日本語訳}

\begin{itemize}
\item Data Driven Decision Making
\item 第2回 - 情報モデリングと意思決定モデリング
\end{itemize}

\end{multicols}

\clearpage

\SlideHeader{Lecture 2: Information and Decision Modeling}{002}{本日の講義目標}

\SlideImage{2}{../Materials/2026/3DM_02_Information-and-Decision-Modeling_SS26.pdf}

\vspace{1.5mm}

\begin{multicols}{2}

\footnotesize

\raggedcolumns

\NoteSection{日本語訳}

\begin{itemize}
\item Business Analytics のプロセス。
\item 情報モデリングと意思決定モデリング。
\item 時間依存旅行時間を持つ車両ルーティング。
\item データが最初に来る。
\item 注意: この講義では多くの手法を使うが、詳細な定義ではなく考え方のみを扱う。
\end{itemize}

\NoteSection{解説}
Real-World は、配送、交通、顧客、車両、作業員など、実際に行動が行われる対象である。ここでは全ての詳細をモデルに入れることはできないため、意思決定に必要な情報だけを選ぶ。\par

Aggregation は、観測データをモデルで使える次元へ変換する処理である。例えばGPS位置列を道路セグメント別の平均速度へ変換すること、住所を顧客ノードへ変換することが該当する。\par

Information Model は、現実を圧縮した情報表現である。旅行時間行列、時間帯別速度パターン、顧客集合、需要分布などが含まれる。\par

Decision Model は、可能な決定、目的関数、制約、評価方法を定義する。TSP/VRPであれば、訪問順序、利用アーク、各顧客を一度訪問する制約、総旅行時間最小化が含まれる。\par

Implementation は、モデルの解を現実の行動へ戻す処理である。例えば『1-3-2』というノード順序を、実際の住所と道路ナビゲーションへ変換する。\par

\NoteSection{試験対策ポイント}
Real-World, Aggregation, Information Model, Decision Model, Implementation の役割とデータ・解の流れを、配送例に対応付けて説明できること。\par

\NoteSection{過去問関連}
\ExamItem{WS23/24 Aufgabe 2 (5点)}{Real-World Problem と Decision Model の関係、および Information Model を介した流れを説明する問題。}{現実世界から直接最適化モデルを作るのではなく、まず観測データを集約して意思決定に必要な情報表現に落とす。Information Model は現実を圧縮した表現で、Decision Model はその表現を使って決定空間、制約、目的関数を定義する。解は Implementation により現実の行動へ変換され、結果観測から仮説・モデルが更新される。}
\ExamItem{SS24 Aufgabe 1a (8点)}{Information Model, Decision Model, Aggregation, Disaggregation/Implementation の機能を説明する問題。}{Aggregation は現実データをモデルの次元へ変換する。Information Model は意思決定に必要な状態・属性・パラメータを整理する。Decision Model は可能な決定、目的、制約、評価方法を定義する。Disaggregation/Implementation はモデル解を現実の作業指示やルートへ戻す。重要なのは、これらが一方向で終わらず、実装後の観測で更新されるループである点である。}
\ExamItem{WS22/23 Aufgabe 3 (6点)}{Business Analytics プロセスの中で OR 的なモデル化・解法がどこにあるかを説明する問題。}{情報モデルはデータ側、意思決定モデルはOR側の橋渡しである。ORは、現実の問題構造を抽象化し、目的関数・制約・決定変数に変換し、解を評価可能にする。}
\ExamItem{WS22/23 Aufgabe 2 (6点)}{Business Analytics を構成する分野と、それらの交差領域を説明する問題。}{Operations Research は意思決定モデル、制約、目的関数、最適化を担う。Statistics はデータの不確実性、分布、推定、予測を扱う。Computer Science/Information Systems はデータ管理、アルゴリズム、実装、情報システムを担う。Business Analytics はこれらを組み合わせ、データから意思決定可能な知見と実行可能な解を作る分野である。}
\ExamItem{WS22/23 Aufgabe 3 (6点)}{Business Analytics の適用プロセスを挙げ、Operations Research に属する工程を示す問題。}{現実問題を観察し、データを取得・集約して情報モデルを作り、そこから意思決定モデルを定義し、解法で解を求め、実装して現実に戻す。ORに特に属するのは、目的関数・制約・決定変数を持つ意思決定モデルの定式化、解空間の探索、最適化・ヒューリスティックによる解の計算である。}
\ExamItem{WS23/24 Aufgabe 1 (6点)}{Business Analytics の三つの構成分野と交差領域を説明する問題。}{三分野を単に列挙するだけでなく、データ分析だけでは意思決定にならず、最適化だけでも現実データを反映できない点を述べる。交差領域では、データからモデルを作り、モデルから現実に実装可能な意思決定を作ることが中心になる。}
\ExamItem{WS23/24 Aufgabe 3 (10点)}{配送車両の業務的ツアープランニングに Business Analytics の構成要素を対応付ける問題。}{Real-world は配送現場、顧客、道路、車両である。Aggregation は住所・道路・走行時間観測をモデル変数へ変換すること。Information model は顧客集合、距離/旅行時間行列、時間帯情報である。Decision model は訪問順序・経路・目的関数・制約を持つVRP/TSPである。Implementation は得られたツアーをドライバーのルート案内として実行することである。}

\NoteSection{確認問題と解答}
\begin{enumerate}
\item \textbf{Information Model と Decision Model の違いを説明せよ。}\\
Information Model は、意思決定に必要な現実情報の表現である。例として顧客、道路、旅行時間行列、時間帯別速度がある。Decision Model は、その情報を用いて決定を選ぶための目的関数、制約、決定変数の体系である。情報モデルが『何を知っているか』を表し、意思決定モデルが『その情報で何をどう決めるか』を表す。
\item \textbf{Aggregation と Implementation の役割を説明せよ。}\\
Aggregation は現実データを抽象化してモデル入力に変換する処理であり、Implementation はモデル解を実際の行動へ具体化する処理である。前者は現実からモデルへの橋、後者はモデルから現実への橋である。
\item \textbf{なぜこのループは一回で終わらないのか。}\\
実装後に現実の結果が観測されるためである。予測より配送が遅れた、交通パターンが違った、制約が現実に合わなかった場合、仮説、情報モデル、意思決定モデルを更新する必要がある。
\end{enumerate}

\end{multicols}

\clearpage

\SlideHeader{Lecture 2: Information and Decision Modeling}{003}{アジェンダ}

\SlideImage{3}{../Materials/2026/3DM_02_Information-and-Decision-Modeling_SS26.pdf}

\vspace{1.5mm}

\begin{multicols}{2}

\footnotesize

\raggedcolumns

\NoteSection{日本語訳}

\begin{itemize}
\item 情報モデリングと意思決定モデリング。
\item ケーススタディ: 時間依存旅行時間を持つ巡回セールスマン問題。
\end{itemize}

\end{multicols}

\clearpage

\SlideHeader{Lecture 2: Information and Decision Modeling}{004}{意思決定}

\SlideImage{4}{../Materials/2026/3DM_02_Information-and-Decision-Modeling_SS26.pdf}

\vspace{1.5mm}

\begin{multicols}{2}

\footnotesize

\raggedcolumns

\NoteSection{日本語訳}

\begin{itemize}
\item Real-World: 行動が実行され、その結果として状態を観測できる。
\item Aggregation: 現実世界のデータを情報モデルの次元へ移す。
\item Information Model: 情報を圧縮された形で描写する。
\item Decision Model: 決定空間を定義し、決定を評価する。
\item Implementation: 意思決定モデルの解を現実世界の行動へ変換する。
\end{itemize}

\NoteSection{解説}
Real-World は、配送、交通、顧客、車両、作業員など、実際に行動が行われる対象である。ここでは全ての詳細をモデルに入れることはできないため、意思決定に必要な情報だけを選ぶ。\par

Aggregation は、観測データをモデルで使える次元へ変換する処理である。例えばGPS位置列を道路セグメント別の平均速度へ変換すること、住所を顧客ノードへ変換することが該当する。\par

Information Model は、現実を圧縮した情報表現である。旅行時間行列、時間帯別速度パターン、顧客集合、需要分布などが含まれる。\par

Decision Model は、可能な決定、目的関数、制約、評価方法を定義する。TSP/VRPであれば、訪問順序、利用アーク、各顧客を一度訪問する制約、総旅行時間最小化が含まれる。\par

Implementation は、モデルの解を現実の行動へ戻す処理である。例えば『1-3-2』というノード順序を、実際の住所と道路ナビゲーションへ変換する。\par

\NoteSection{試験対策ポイント}
Real-World, Aggregation, Information Model, Decision Model, Implementation の役割とデータ・解の流れを、配送例に対応付けて説明できること。\par

\NoteSection{過去問関連}
\ExamItem{WS23/24 Aufgabe 2 (5点)}{Real-World Problem と Decision Model の関係、および Information Model を介した流れを説明する問題。}{現実世界から直接最適化モデルを作るのではなく、まず観測データを集約して意思決定に必要な情報表現に落とす。Information Model は現実を圧縮した表現で、Decision Model はその表現を使って決定空間、制約、目的関数を定義する。解は Implementation により現実の行動へ変換され、結果観測から仮説・モデルが更新される。}
\ExamItem{SS24 Aufgabe 1a (8点)}{Information Model, Decision Model, Aggregation, Disaggregation/Implementation の機能を説明する問題。}{Aggregation は現実データをモデルの次元へ変換する。Information Model は意思決定に必要な状態・属性・パラメータを整理する。Decision Model は可能な決定、目的、制約、評価方法を定義する。Disaggregation/Implementation はモデル解を現実の作業指示やルートへ戻す。重要なのは、これらが一方向で終わらず、実装後の観測で更新されるループである点である。}
\ExamItem{WS22/23 Aufgabe 3 (6点)}{Business Analytics プロセスの中で OR 的なモデル化・解法がどこにあるかを説明する問題。}{情報モデルはデータ側、意思決定モデルはOR側の橋渡しである。ORは、現実の問題構造を抽象化し、目的関数・制約・決定変数に変換し、解を評価可能にする。}
\ExamItem{WS22/23 Aufgabe 2 (6点)}{Business Analytics を構成する分野と、それらの交差領域を説明する問題。}{Operations Research は意思決定モデル、制約、目的関数、最適化を担う。Statistics はデータの不確実性、分布、推定、予測を扱う。Computer Science/Information Systems はデータ管理、アルゴリズム、実装、情報システムを担う。Business Analytics はこれらを組み合わせ、データから意思決定可能な知見と実行可能な解を作る分野である。}
\ExamItem{WS22/23 Aufgabe 3 (6点)}{Business Analytics の適用プロセスを挙げ、Operations Research に属する工程を示す問題。}{現実問題を観察し、データを取得・集約して情報モデルを作り、そこから意思決定モデルを定義し、解法で解を求め、実装して現実に戻す。ORに特に属するのは、目的関数・制約・決定変数を持つ意思決定モデルの定式化、解空間の探索、最適化・ヒューリスティックによる解の計算である。}
\ExamItem{WS23/24 Aufgabe 1 (6点)}{Business Analytics の三つの構成分野と交差領域を説明する問題。}{三分野を単に列挙するだけでなく、データ分析だけでは意思決定にならず、最適化だけでも現実データを反映できない点を述べる。交差領域では、データからモデルを作り、モデルから現実に実装可能な意思決定を作ることが中心になる。}
\ExamItem{WS23/24 Aufgabe 3 (10点)}{配送車両の業務的ツアープランニングに Business Analytics の構成要素を対応付ける問題。}{Real-world は配送現場、顧客、道路、車両である。Aggregation は住所・道路・走行時間観測をモデル変数へ変換すること。Information model は顧客集合、距離/旅行時間行列、時間帯情報である。Decision model は訪問順序・経路・目的関数・制約を持つVRP/TSPである。Implementation は得られたツアーをドライバーのルート案内として実行することである。}

\NoteSection{確認問題と解答}
\begin{enumerate}
\item \textbf{Information Model と Decision Model の違いを説明せよ。}\\
Information Model は、意思決定に必要な現実情報の表現である。例として顧客、道路、旅行時間行列、時間帯別速度がある。Decision Model は、その情報を用いて決定を選ぶための目的関数、制約、決定変数の体系である。情報モデルが『何を知っているか』を表し、意思決定モデルが『その情報で何をどう決めるか』を表す。
\item \textbf{Aggregation と Implementation の役割を説明せよ。}\\
Aggregation は現実データを抽象化してモデル入力に変換する処理であり、Implementation はモデル解を実際の行動へ具体化する処理である。前者は現実からモデルへの橋、後者はモデルから現実への橋である。
\item \textbf{なぜこのループは一回で終わらないのか。}\\
実装後に現実の結果が観測されるためである。予測より配送が遅れた、交通パターンが違った、制約が現実に合わなかった場合、仮説、情報モデル、意思決定モデルを更新する必要がある。
\end{enumerate}

\end{multicols}

\clearpage

\SlideHeader{Lecture 2: Information and Decision Modeling}{005}{情報モデルと意思決定モデルの選択}

\SlideImage{5}{../Materials/2026/3DM_02_Information-and-Decision-Modeling_SS26.pdf}

\vspace{1.5mm}

\begin{multicols}{2}

\footnotesize

\raggedcolumns

\NoteSection{日本語訳}

\begin{itemize}
\item 情報モデルと意思決定モデルは初期仮説に基づいて設計される。
\item 仮説は現実問題の観測行動とデータに基づく。これは descriptive analytics に対応する。
\item 仮説は必ずしも正しいとは限らない。
\item 仮説の品質は、意思決定モデルの解を実装することで評価できる。
\item 観測結果により仮説が更新され、その結果として情報モデルと意思決定モデルも更新される。
\end{itemize}

\NoteSection{解説}
Real-World は、配送、交通、顧客、車両、作業員など、実際に行動が行われる対象である。ここでは全ての詳細をモデルに入れることはできないため、意思決定に必要な情報だけを選ぶ。\par

Aggregation は、観測データをモデルで使える次元へ変換する処理である。例えばGPS位置列を道路セグメント別の平均速度へ変換すること、住所を顧客ノードへ変換することが該当する。\par

Information Model は、現実を圧縮した情報表現である。旅行時間行列、時間帯別速度パターン、顧客集合、需要分布などが含まれる。\par

Decision Model は、可能な決定、目的関数、制約、評価方法を定義する。TSP/VRPであれば、訪問順序、利用アーク、各顧客を一度訪問する制約、総旅行時間最小化が含まれる。\par

Implementation は、モデルの解を現実の行動へ戻す処理である。例えば『1-3-2』というノード順序を、実際の住所と道路ナビゲーションへ変換する。\par

\NoteSection{試験対策ポイント}
Real-World, Aggregation, Information Model, Decision Model, Implementation の役割とデータ・解の流れを、配送例に対応付けて説明できること。\par

\NoteSection{過去問関連}
\ExamItem{WS23/24 Aufgabe 2 (5点)}{Real-World Problem と Decision Model の関係、および Information Model を介した流れを説明する問題。}{現実世界から直接最適化モデルを作るのではなく、まず観測データを集約して意思決定に必要な情報表現に落とす。Information Model は現実を圧縮した表現で、Decision Model はその表現を使って決定空間、制約、目的関数を定義する。解は Implementation により現実の行動へ変換され、結果観測から仮説・モデルが更新される。}
\ExamItem{SS24 Aufgabe 1a (8点)}{Information Model, Decision Model, Aggregation, Disaggregation/Implementation の機能を説明する問題。}{Aggregation は現実データをモデルの次元へ変換する。Information Model は意思決定に必要な状態・属性・パラメータを整理する。Decision Model は可能な決定、目的、制約、評価方法を定義する。Disaggregation/Implementation はモデル解を現実の作業指示やルートへ戻す。重要なのは、これらが一方向で終わらず、実装後の観測で更新されるループである点である。}
\ExamItem{WS22/23 Aufgabe 3 (6点)}{Business Analytics プロセスの中で OR 的なモデル化・解法がどこにあるかを説明する問題。}{情報モデルはデータ側、意思決定モデルはOR側の橋渡しである。ORは、現実の問題構造を抽象化し、目的関数・制約・決定変数に変換し、解を評価可能にする。}
\ExamItem{WS22/23 Aufgabe 2 (6点)}{Business Analytics を構成する分野と、それらの交差領域を説明する問題。}{Operations Research は意思決定モデル、制約、目的関数、最適化を担う。Statistics はデータの不確実性、分布、推定、予測を扱う。Computer Science/Information Systems はデータ管理、アルゴリズム、実装、情報システムを担う。Business Analytics はこれらを組み合わせ、データから意思決定可能な知見と実行可能な解を作る分野である。}
\ExamItem{WS22/23 Aufgabe 3 (6点)}{Business Analytics の適用プロセスを挙げ、Operations Research に属する工程を示す問題。}{現実問題を観察し、データを取得・集約して情報モデルを作り、そこから意思決定モデルを定義し、解法で解を求め、実装して現実に戻す。ORに特に属するのは、目的関数・制約・決定変数を持つ意思決定モデルの定式化、解空間の探索、最適化・ヒューリスティックによる解の計算である。}
\ExamItem{WS23/24 Aufgabe 1 (6点)}{Business Analytics の三つの構成分野と交差領域を説明する問題。}{三分野を単に列挙するだけでなく、データ分析だけでは意思決定にならず、最適化だけでも現実データを反映できない点を述べる。交差領域では、データからモデルを作り、モデルから現実に実装可能な意思決定を作ることが中心になる。}
\ExamItem{WS23/24 Aufgabe 3 (10点)}{配送車両の業務的ツアープランニングに Business Analytics の構成要素を対応付ける問題。}{Real-world は配送現場、顧客、道路、車両である。Aggregation は住所・道路・走行時間観測をモデル変数へ変換すること。Information model は顧客集合、距離/旅行時間行列、時間帯情報である。Decision model は訪問順序・経路・目的関数・制約を持つVRP/TSPである。Implementation は得られたツアーをドライバーのルート案内として実行することである。}

\NoteSection{確認問題と解答}
\begin{enumerate}
\item \textbf{Information Model と Decision Model の違いを説明せよ。}\\
Information Model は、意思決定に必要な現実情報の表現である。例として顧客、道路、旅行時間行列、時間帯別速度がある。Decision Model は、その情報を用いて決定を選ぶための目的関数、制約、決定変数の体系である。情報モデルが『何を知っているか』を表し、意思決定モデルが『その情報で何をどう決めるか』を表す。
\item \textbf{Aggregation と Implementation の役割を説明せよ。}\\
Aggregation は現実データを抽象化してモデル入力に変換する処理であり、Implementation はモデル解を実際の行動へ具体化する処理である。前者は現実からモデルへの橋、後者はモデルから現実への橋である。
\item \textbf{なぜこのループは一回で終わらないのか。}\\
実装後に現実の結果が観測されるためである。予測より配送が遅れた、交通パターンが違った、制約が現実に合わなかった場合、仮説、情報モデル、意思決定モデルを更新する必要がある。
\end{enumerate}

\end{multicols}

\clearpage

\SlideHeader{Lecture 2: Information and Decision Modeling}{006}{IDA: Appearance から Structureへ}

\SlideImage{6}{../Materials/2026/3DM_02_Information-and-Decision-Modeling_SS26.pdf}

\vspace{1.5mm}

\begin{multicols}{2}

\footnotesize

\raggedcolumns

\NoteSection{日本語訳}

\begin{itemize}
\item 現実システムからのデータは、その見かけ appearance を反映するだけで、構造は直接見えにくい。
\item appearance は不完全、誤り、例外を含む可能性がある。
\item システム構造はデータモデル構築を助ける。
\item 記録データは一般的なシステムappearanceの少数サンプルにすぎない。
\item 記録データは、クリーニング、修正、補完、検証、絞り込みによりターゲットデータセットへ変換される。
\item データマイニングはターゲットデータからモデルを計算する。
\item データマイニング概念は、データ内の統計的相関に関する仮説を実装する。
\item モデルまたはパターンはシステム構造に関する証拠を与える。
\end{itemize}

\NoteSection{解説}
IDA は、観測されたデータの appearance から、背後にある system structure を推測する方向である。データのクリーニング、補完、検証、データマイニング、パターン抽出が含まれる。\par

PMT/OR は、問題構造から意思決定モデルを作り、解を計算して現実に実装する方向である。目的、決定空間、制約、ルーティング・スケジューリング・割当のようなモデル構造が重要になる。\par

統一的に見ると、Business Analytics はデータから情報モデルを作る流れと、問題構造から意思決定モデルを作る流れを接続する。DDDMで重要なのは、どちらか一方ではなく、情報モデルと意思決定モデルが互いに依存する点である。\par

\NoteSection{試験対策ポイント}
IDAはデータから構造へ、OR/PMTは構造から解へ向かう。この二つがBusiness Analyticsで接続されることを図なしで説明できること。\par

\NoteSection{過去問関連}
\ExamItem{WS22/23 Aufgabe 2 (6点)}{Business Analytics を構成する分野と、それらの交差領域を説明する問題。}{Operations Research は意思決定モデル、制約、目的関数、最適化を担う。Statistics はデータの不確実性、分布、推定、予測を扱う。Computer Science/Information Systems はデータ管理、アルゴリズム、実装、情報システムを担う。Business Analytics はこれらを組み合わせ、データから意思決定可能な知見と実行可能な解を作る分野である。}
\ExamItem{WS22/23 Aufgabe 3 (6点)}{Business Analytics の適用プロセスを挙げ、Operations Research に属する工程を示す問題。}{現実問題を観察し、データを取得・集約して情報モデルを作り、そこから意思決定モデルを定義し、解法で解を求め、実装して現実に戻す。ORに特に属するのは、目的関数・制約・決定変数を持つ意思決定モデルの定式化、解空間の探索、最適化・ヒューリスティックによる解の計算である。}
\ExamItem{WS23/24 Aufgabe 1 (6点)}{Business Analytics の三つの構成分野と交差領域を説明する問題。}{三分野を単に列挙するだけでなく、データ分析だけでは意思決定にならず、最適化だけでも現実データを反映できない点を述べる。交差領域では、データからモデルを作り、モデルから現実に実装可能な意思決定を作ることが中心になる。}
\ExamItem{WS23/24 Aufgabe 3 (10点)}{配送車両の業務的ツアープランニングに Business Analytics の構成要素を対応付ける問題。}{Real-world は配送現場、顧客、道路、車両である。Aggregation は住所・道路・走行時間観測をモデル変数へ変換すること。Information model は顧客集合、距離/旅行時間行列、時間帯情報である。Decision model は訪問順序・経路・目的関数・制約を持つVRP/TSPである。Implementation は得られたツアーをドライバーのルート案内として実行することである。}
\ExamItem{WS23/24 Aufgabe 2 (5点)}{Real-World Problem と Decision Model の関係、および Information Model を介した流れを説明する問題。}{現実世界から直接最適化モデルを作るのではなく、まず観測データを集約して意思決定に必要な情報表現に落とす。Information Model は現実を圧縮した表現で、Decision Model はその表現を使って決定空間、制約、目的関数を定義する。解は Implementation により現実の行動へ変換され、結果観測から仮説・モデルが更新される。}
\ExamItem{SS24 Aufgabe 1a (8点)}{Information Model, Decision Model, Aggregation, Disaggregation/Implementation の機能を説明する問題。}{Aggregation は現実データをモデルの次元へ変換する。Information Model は意思決定に必要な状態・属性・パラメータを整理する。Decision Model は可能な決定、目的、制約、評価方法を定義する。Disaggregation/Implementation はモデル解を現実の作業指示やルートへ戻す。重要なのは、これらが一方向で終わらず、実装後の観測で更新されるループである点である。}
\ExamItem{WS22/23 Aufgabe 3 (6点)}{Business Analytics プロセスの中で OR 的なモデル化・解法がどこにあるかを説明する問題。}{情報モデルはデータ側、意思決定モデルはOR側の橋渡しである。ORは、現実の問題構造を抽象化し、目的関数・制約・決定変数に変換し、解を評価可能にする。}

\NoteSection{確認問題と解答}
\begin{enumerate}
\item \textbf{IDA と OR/PMT の方向性の違いを説明せよ。}\\
IDAは観測データから構造を見つける。つまり appearance から structure へ向かう。OR/PMTは問題構造を仮説として選び、意思決定モデルと解を作る。つまり structure から solution/appearance へ向かう。Business Analyticsでは両者を循環的に使う。
\item \textbf{なぜデータはそのまま構造を示さないのか。}\\
データは現実の一部の観測であり、不完全、誤り、例外、ノイズを含む。例えばGoogle Mapsが渋滞を示しても、その原因や構造までは直接示さない。したがってデータ処理と仮説に基づくモデル化が必要になる。
\end{enumerate}

\end{multicols}

\clearpage

\SlideHeader{Lecture 2: Information and Decision Modeling}{007}{PMT: Structure から Appearanceへ}

\SlideImage{7}{../Materials/2026/3DM_02_Information-and-Decision-Modeling_SS26.pdf}

\vspace{1.5mm}

\begin{multicols}{2}

\footnotesize

\raggedcolumns

\NoteSection{日本語訳}

\begin{itemize}
\item 現実の意思決定問題は、数理モデリングの要件を満たすよう単純化される。
\item 仮説により重要な問題構造の特徴が選択される。
\item データが利用可能な特徴だけが考慮される。
\item 重要な特徴は目的、決定空間、制約である。
\item 意思決定モデルの一般構造が、問題から意思決定モデルへの道を橋渡しする。
\item 意思決定モデル構造の例は、割当、ルーティング、スケジューリングである。
\item 得られた解は問題の一つの最適な appearance を描く。
\end{itemize}

\NoteSection{解説}
IDA は、観測されたデータの appearance から、背後にある system structure を推測する方向である。データのクリーニング、補完、検証、データマイニング、パターン抽出が含まれる。\par

PMT/OR は、問題構造から意思決定モデルを作り、解を計算して現実に実装する方向である。目的、決定空間、制約、ルーティング・スケジューリング・割当のようなモデル構造が重要になる。\par

統一的に見ると、Business Analytics はデータから情報モデルを作る流れと、問題構造から意思決定モデルを作る流れを接続する。DDDMで重要なのは、どちらか一方ではなく、情報モデルと意思決定モデルが互いに依存する点である。\par

\NoteSection{試験対策ポイント}
IDAはデータから構造へ、OR/PMTは構造から解へ向かう。この二つがBusiness Analyticsで接続されることを図なしで説明できること。\par

\NoteSection{過去問関連}
\ExamItem{WS22/23 Aufgabe 2 (6点)}{Business Analytics を構成する分野と、それらの交差領域を説明する問題。}{Operations Research は意思決定モデル、制約、目的関数、最適化を担う。Statistics はデータの不確実性、分布、推定、予測を扱う。Computer Science/Information Systems はデータ管理、アルゴリズム、実装、情報システムを担う。Business Analytics はこれらを組み合わせ、データから意思決定可能な知見と実行可能な解を作る分野である。}
\ExamItem{WS22/23 Aufgabe 3 (6点)}{Business Analytics の適用プロセスを挙げ、Operations Research に属する工程を示す問題。}{現実問題を観察し、データを取得・集約して情報モデルを作り、そこから意思決定モデルを定義し、解法で解を求め、実装して現実に戻す。ORに特に属するのは、目的関数・制約・決定変数を持つ意思決定モデルの定式化、解空間の探索、最適化・ヒューリスティックによる解の計算である。}
\ExamItem{WS23/24 Aufgabe 1 (6点)}{Business Analytics の三つの構成分野と交差領域を説明する問題。}{三分野を単に列挙するだけでなく、データ分析だけでは意思決定にならず、最適化だけでも現実データを反映できない点を述べる。交差領域では、データからモデルを作り、モデルから現実に実装可能な意思決定を作ることが中心になる。}
\ExamItem{WS23/24 Aufgabe 3 (10点)}{配送車両の業務的ツアープランニングに Business Analytics の構成要素を対応付ける問題。}{Real-world は配送現場、顧客、道路、車両である。Aggregation は住所・道路・走行時間観測をモデル変数へ変換すること。Information model は顧客集合、距離/旅行時間行列、時間帯情報である。Decision model は訪問順序・経路・目的関数・制約を持つVRP/TSPである。Implementation は得られたツアーをドライバーのルート案内として実行することである。}
\ExamItem{WS23/24 Aufgabe 2 (5点)}{Real-World Problem と Decision Model の関係、および Information Model を介した流れを説明する問題。}{現実世界から直接最適化モデルを作るのではなく、まず観測データを集約して意思決定に必要な情報表現に落とす。Information Model は現実を圧縮した表現で、Decision Model はその表現を使って決定空間、制約、目的関数を定義する。解は Implementation により現実の行動へ変換され、結果観測から仮説・モデルが更新される。}
\ExamItem{SS24 Aufgabe 1a (8点)}{Information Model, Decision Model, Aggregation, Disaggregation/Implementation の機能を説明する問題。}{Aggregation は現実データをモデルの次元へ変換する。Information Model は意思決定に必要な状態・属性・パラメータを整理する。Decision Model は可能な決定、目的、制約、評価方法を定義する。Disaggregation/Implementation はモデル解を現実の作業指示やルートへ戻す。重要なのは、これらが一方向で終わらず、実装後の観測で更新されるループである点である。}
\ExamItem{WS22/23 Aufgabe 3 (6点)}{Business Analytics プロセスの中で OR 的なモデル化・解法がどこにあるかを説明する問題。}{情報モデルはデータ側、意思決定モデルはOR側の橋渡しである。ORは、現実の問題構造を抽象化し、目的関数・制約・決定変数に変換し、解を評価可能にする。}

\NoteSection{確認問題と解答}
\begin{enumerate}
\item \textbf{IDA と OR/PMT の方向性の違いを説明せよ。}\\
IDAは観測データから構造を見つける。つまり appearance から structure へ向かう。OR/PMTは問題構造を仮説として選び、意思決定モデルと解を作る。つまり structure から solution/appearance へ向かう。Business Analyticsでは両者を循環的に使う。
\item \textbf{なぜデータはそのまま構造を示さないのか。}\\
データは現実の一部の観測であり、不完全、誤り、例外、ノイズを含む。例えばGoogle Mapsが渋滞を示しても、その原因や構造までは直接示さない。したがってデータ処理と仮説に基づくモデル化が必要になる。
\end{enumerate}

\end{multicols}

\clearpage

\SlideHeader{Lecture 2: Information and Decision Modeling}{008}{IDAとPMTの統一的見方}

\SlideImage{8}{../Materials/2026/3DM_02_Information-and-Decision-Modeling_SS26.pdf}

\vspace{1.5mm}

\begin{multicols}{2}

\footnotesize

\raggedcolumns

\NoteSection{日本語訳}

\begin{itemize}
\item 記録データからターゲットデータへ。
\item ターゲットデータから情報構造と情報モデル/パターンへ。
\item 問題構造から意思決定モデル構造と意思決定モデルへ。
\item 仮説が、データ側と意思決定側の両方を接続する。
\item 解は現実で実装され、再び観測につながる。
\end{itemize}

\NoteSection{解説}
IDA は、観測されたデータの appearance から、背後にある system structure を推測する方向である。データのクリーニング、補完、検証、データマイニング、パターン抽出が含まれる。\par

PMT/OR は、問題構造から意思決定モデルを作り、解を計算して現実に実装する方向である。目的、決定空間、制約、ルーティング・スケジューリング・割当のようなモデル構造が重要になる。\par

統一的に見ると、Business Analytics はデータから情報モデルを作る流れと、問題構造から意思決定モデルを作る流れを接続する。DDDMで重要なのは、どちらか一方ではなく、情報モデルと意思決定モデルが互いに依存する点である。\par

\NoteSection{試験対策ポイント}
IDAはデータから構造へ、OR/PMTは構造から解へ向かう。この二つがBusiness Analyticsで接続されることを図なしで説明できること。\par

\NoteSection{過去問関連}
\ExamItem{WS22/23 Aufgabe 2 (6点)}{Business Analytics を構成する分野と、それらの交差領域を説明する問題。}{Operations Research は意思決定モデル、制約、目的関数、最適化を担う。Statistics はデータの不確実性、分布、推定、予測を扱う。Computer Science/Information Systems はデータ管理、アルゴリズム、実装、情報システムを担う。Business Analytics はこれらを組み合わせ、データから意思決定可能な知見と実行可能な解を作る分野である。}
\ExamItem{WS22/23 Aufgabe 3 (6点)}{Business Analytics の適用プロセスを挙げ、Operations Research に属する工程を示す問題。}{現実問題を観察し、データを取得・集約して情報モデルを作り、そこから意思決定モデルを定義し、解法で解を求め、実装して現実に戻す。ORに特に属するのは、目的関数・制約・決定変数を持つ意思決定モデルの定式化、解空間の探索、最適化・ヒューリスティックによる解の計算である。}
\ExamItem{WS23/24 Aufgabe 1 (6点)}{Business Analytics の三つの構成分野と交差領域を説明する問題。}{三分野を単に列挙するだけでなく、データ分析だけでは意思決定にならず、最適化だけでも現実データを反映できない点を述べる。交差領域では、データからモデルを作り、モデルから現実に実装可能な意思決定を作ることが中心になる。}
\ExamItem{WS23/24 Aufgabe 3 (10点)}{配送車両の業務的ツアープランニングに Business Analytics の構成要素を対応付ける問題。}{Real-world は配送現場、顧客、道路、車両である。Aggregation は住所・道路・走行時間観測をモデル変数へ変換すること。Information model は顧客集合、距離/旅行時間行列、時間帯情報である。Decision model は訪問順序・経路・目的関数・制約を持つVRP/TSPである。Implementation は得られたツアーをドライバーのルート案内として実行することである。}
\ExamItem{WS23/24 Aufgabe 2 (5点)}{Real-World Problem と Decision Model の関係、および Information Model を介した流れを説明する問題。}{現実世界から直接最適化モデルを作るのではなく、まず観測データを集約して意思決定に必要な情報表現に落とす。Information Model は現実を圧縮した表現で、Decision Model はその表現を使って決定空間、制約、目的関数を定義する。解は Implementation により現実の行動へ変換され、結果観測から仮説・モデルが更新される。}
\ExamItem{SS24 Aufgabe 1a (8点)}{Information Model, Decision Model, Aggregation, Disaggregation/Implementation の機能を説明する問題。}{Aggregation は現実データをモデルの次元へ変換する。Information Model は意思決定に必要な状態・属性・パラメータを整理する。Decision Model は可能な決定、目的、制約、評価方法を定義する。Disaggregation/Implementation はモデル解を現実の作業指示やルートへ戻す。重要なのは、これらが一方向で終わらず、実装後の観測で更新されるループである点である。}
\ExamItem{WS22/23 Aufgabe 3 (6点)}{Business Analytics プロセスの中で OR 的なモデル化・解法がどこにあるかを説明する問題。}{情報モデルはデータ側、意思決定モデルはOR側の橋渡しである。ORは、現実の問題構造を抽象化し、目的関数・制約・決定変数に変換し、解を評価可能にする。}

\NoteSection{確認問題と解答}
\begin{enumerate}
\item \textbf{IDA と OR/PMT の方向性の違いを説明せよ。}\\
IDAは観測データから構造を見つける。つまり appearance から structure へ向かう。OR/PMTは問題構造を仮説として選び、意思決定モデルと解を作る。つまり structure から solution/appearance へ向かう。Business Analyticsでは両者を循環的に使う。
\item \textbf{なぜデータはそのまま構造を示さないのか。}\\
データは現実の一部の観測であり、不完全、誤り、例外、ノイズを含む。例えばGoogle Mapsが渋滞を示しても、その原因や構造までは直接示さない。したがってデータ処理と仮説に基づくモデル化が必要になる。
\end{enumerate}

\end{multicols}

\clearpage

\SlideHeader{Lecture 2: Information and Decision Modeling}{009}{例: 巡回セールスマン問題}

\SlideImage{9}{../Materials/2026/3DM_02_Information-and-Decision-Modeling_SS26.pdf}

\vspace{1.5mm}

\begin{multicols}{2}

\footnotesize

\raggedcolumns

\NoteSection{日本語訳}

\begin{itemize}
\item 仮説: 旅行時間は一定である。
\item 情報モデル: 旅行時間行列。
\item 意思決定モデル: 時間的考慮を持たない混合整数計画。
\item 実装すると予想外の結果が生じる。通常、計画より時間がかかる。
\item 仮説を、時間依存旅行時間を捉えるように調整する。
\begin{itemize}
\item 情報モデル: 複数の旅行時間行列と遷移に関する考慮。
\item 意思決定モデル: 時刻を考慮した混合整数計画。決定変数は x\_ij(t) のようになる。
\end{itemize}
\end{itemize}

\NoteSection{解説}
TSP は、顧客をノード、移動可能性をアーク、移動時間または距離を重みとして表し、全顧客を一度ずつ訪問して出発点に戻る経路を求める問題である。\par

情報モデルは、道路ネットワークや旅行速度を、顧客間の旅行時間行列へ変換する部分である。意思決定モデルは、どのアークを使うかを二値変数で表し、各顧客の入次数・出次数、部分巡回除去、目的関数を定式化する部分である。\par

静的TSPでは旅行時間は定数と仮定する。しかし都市配送では出発時刻により旅行時間が変わるため、この仮説が実装後の遅延や残業として破綻することがある。\par

\NoteSection{試験対策ポイント}
TSPの情報モデルと意思決定モデルを分けて説明し、静的旅行時間仮説が実装後に崩れる理由を説明できること。\par

\NoteSection{過去問関連}
\ExamItem{WS22/23 Aufgabe 1 (7点)}{TSP の数理最適化モデルと、時間次元を持つTSPで何が変わるかを問う問題。}{標準TSPでは、顧客ノード、アーク、旅行時間または距離、二値変数 x\_ij、各顧客を一度訪問する制約、部分巡回除去制約、総旅行時間最小化を定義する。時間依存TSPでは旅行時間が定数 d\_ij ではなく出発時刻 t に依存する d\_ij(t) となり、あるアーク選択が後続の到着時刻と後続アークの旅行時間を変えるため、静的モデルより決定変数・制約・解法が複雑になる。}
\ExamItem{SS24 Aufgabe 1b (4点)}{旅行時間の情報モデルの例と、常に最も詳細なモデルを選ばない理由を問う問題。}{例として、単一平均旅行時間行列、時間帯別行列、区間旅行時間、確率分布モデルがある。詳細なモデルほど現実に近い可能性はあるが、データ量、信頼性、計算時間、解釈可能性、意思決定モデルの複雑性が悪化する。したがって目的に十分な粒度を選ぶ。}
\ExamItem{WS23/24 Aufgabe 2 (5点)}{Real-World Problem と Decision Model の関係、および Information Model を介した流れを説明する問題。}{現実世界から直接最適化モデルを作るのではなく、まず観測データを集約して意思決定に必要な情報表現に落とす。Information Model は現実を圧縮した表現で、Decision Model はその表現を使って決定空間、制約、目的関数を定義する。解は Implementation により現実の行動へ変換され、結果観測から仮説・モデルが更新される。}
\ExamItem{SS24 Aufgabe 1a (8点)}{Information Model, Decision Model, Aggregation, Disaggregation/Implementation の機能を説明する問題。}{Aggregation は現実データをモデルの次元へ変換する。Information Model は意思決定に必要な状態・属性・パラメータを整理する。Decision Model は可能な決定、目的、制約、評価方法を定義する。Disaggregation/Implementation はモデル解を現実の作業指示やルートへ戻す。重要なのは、これらが一方向で終わらず、実装後の観測で更新されるループである点である。}
\ExamItem{WS22/23 Aufgabe 3 (6点)}{Business Analytics プロセスの中で OR 的なモデル化・解法がどこにあるかを説明する問題。}{情報モデルはデータ側、意思決定モデルはOR側の橋渡しである。ORは、現実の問題構造を抽象化し、目的関数・制約・決定変数に変換し、解を評価可能にする。}

\NoteSection{確認問題と解答}
\begin{enumerate}
\item \textbf{TSPの情報モデルと意思決定モデルを分けて説明せよ。}\\
情報モデルは顧客、デポ、顧客間旅行時間行列である。意思決定モデルは、アーク選択変数 x\_ij、各顧客を一度訪問する制約、部分巡回除去制約、総旅行時間最小化の目的関数からなる。
\item \textbf{静的TSPで『旅行時間は一定』という仮説がなぜ問題になるか。}\\
都市交通では朝夕のラッシュ、場所ごとの混雑、曜日差により旅行時間が変わる。旅行時間を一定とすると、モデル上は実行可能で短いルートでも、現実では時間超過や顧客不在が起きる可能性がある。
\end{enumerate}

\end{multicols}

\clearpage

\SlideHeader{Lecture 2: Information and Decision Modeling}{010}{アジェンダ}

\SlideImage{10}{../Materials/2026/3DM_02_Information-and-Decision-Modeling_SS26.pdf}

\vspace{1.5mm}

\begin{multicols}{2}

\footnotesize

\raggedcolumns

\NoteSection{日本語訳}

\begin{itemize}
\item 情報モデリングと意思決定モデリング。
\item ケーススタディ: 時間依存旅行時間を持つ巡回セールスマン問題。
\end{itemize}

\end{multicols}

\clearpage

\SlideHeader{Lecture 2: Information and Decision Modeling}{011}{動機}

\SlideImage{11}{../Materials/2026/3DM_02_Information-and-Decision-Modeling_SS26.pdf}

\vspace{1.5mm}

\begin{multicols}{2}

\footnotesize

\raggedcolumns

\NoteSection{日本語訳}

\begin{itemize}
\item サービスルーティング。
\item 高価な走行時間を節約する。
\item 時間的コミットメント
\begin{itemize}
\item 顧客サービス経験を向上させる。
\item 顧客の在宅・利用可能性を確保する。
\end{itemize}
\item どちらも旅行時間の正確な考慮を必要とする。
\end{itemize}

\NoteSection{解説}
時間依存旅行時間では、顧客 i から j への旅行時間を定数ではなく出発時刻 t の関数として扱う。記号的には d\_ij(t) または tau\_ij(t) と考える。\par

同じ二地点間でも、出発時刻が異なれば最短経路も旅行時間も変わる。したがって、訪問順序の変更は単にアークの足し算を変えるだけでなく、後続の到着時刻と全体の旅行時間を変える。\par

このため、情報モデルは時間帯別行列や時間依存リンク旅行時間を必要とし、意思決定モデルは時間の次元を含む変数・制約を必要とする。\par

\NoteSection{試験対策ポイント}
旅行時間を d\_ij(t) として捉え、出発時刻が経路・順序・後続旅行時間に与える影響を説明できること。\par

\NoteSection{過去問関連}
\ExamItem{WS22/23 Aufgabe 1 (7点)}{TSP の数理最適化モデルと、時間次元を持つTSPで何が変わるかを問う問題。}{標準TSPでは、顧客ノード、アーク、旅行時間または距離、二値変数 x\_ij、各顧客を一度訪問する制約、部分巡回除去制約、総旅行時間最小化を定義する。時間依存TSPでは旅行時間が定数 d\_ij ではなく出発時刻 t に依存する d\_ij(t) となり、あるアーク選択が後続の到着時刻と後続アークの旅行時間を変えるため、静的モデルより決定変数・制約・解法が複雑になる。}
\ExamItem{SS24 Aufgabe 1b (4点)}{旅行時間の情報モデルの例と、常に最も詳細なモデルを選ばない理由を問う問題。}{例として、単一平均旅行時間行列、時間帯別行列、区間旅行時間、確率分布モデルがある。詳細なモデルほど現実に近い可能性はあるが、データ量、信頼性、計算時間、解釈可能性、意思決定モデルの複雑性が悪化する。したがって目的に十分な粒度を選ぶ。}
\ExamItem{WS22/23 Aufgabe 4 (8点)}{過去の交通観測から時間帯別平均旅行時間を構築する手順を説明する問題。}{まずFCDなどから位置・時刻・速度を収集する。位置を道路セグメントにマッチングし、セグメントごとの観測速度を得る。外れ値やエラーを除去し、時間帯 w とセグメント l ごとに観測を集約する。平均速度または旅行時間を計算し、観測不足の部分は類似セグメントや時間帯で一般化する。最後に時間帯別の旅行時間行列を作り、TSP/VRPの情報モデルとして使う。}
\ExamItem{SS24 Aufgabe 1b (4点)}{旅行時間情報モデルの選択と複雑性のトレードオフを問う問題。}{平均値モデルは簡単だがラッシュアワーを無視しやすい。時間帯別モデルは実際の交通変動を反映するが、観測数が少ないセルでは信頼性が低くなる。したがってデータ量と意思決定への効果を見ながら集約・一般化する必要がある。}

\NoteSection{確認問題と解答}
\begin{enumerate}
\item \textbf{時間依存旅行時間モデルでは、なぜ一つの旅行時間行列では不十分か。}\\
一つの行列では全時間帯で同じ旅行時間を仮定することになる。実際には道路セグメントごとに混雑時間が異なり、10時と17時では同じ顧客間の最短経路や所要時間が異なる。そのため時間帯別または連続時刻の情報モデルが必要になる。
\item \textbf{時間依存性がヒューリスティックを難しくする理由を説明せよ。}\\
ヒューリスティックの局所的な選択が、後続の出発時刻を変え、後続の旅行時間関数の評価値を変えるからである。静的な追加距離だけでは全体影響を評価できない。
\end{enumerate}

\end{multicols}

\clearpage

\SlideHeader{Lecture 2: Information and Decision Modeling}{012}{時間依存旅行時間}

\SlideImage{12}{../Materials/2026/3DM_02_Information-and-Decision-Modeling_SS26.pdf}

\vspace{1.5mm}

\begin{multicols}{2}

\footnotesize

\raggedcolumns

\NoteSection{日本語訳}

\begin{itemize}
\item 交通は時間と空間により変化する。
\item 道路セグメントごとに旅行時間が異なる。
\item 顧客間の最短経路が異なる。
\item ツアー開始時刻により顧客の訪問順序が異なる。
\end{itemize}

\NoteSection{解説}
時間依存旅行時間では、顧客 i から j への旅行時間を定数ではなく出発時刻 t の関数として扱う。記号的には d\_ij(t) または tau\_ij(t) と考える。\par

同じ二地点間でも、出発時刻が異なれば最短経路も旅行時間も変わる。したがって、訪問順序の変更は単にアークの足し算を変えるだけでなく、後続の到着時刻と全体の旅行時間を変える。\par

このため、情報モデルは時間帯別行列や時間依存リンク旅行時間を必要とし、意思決定モデルは時間の次元を含む変数・制約を必要とする。\par

\NoteSection{試験対策ポイント}
旅行時間を d\_ij(t) として捉え、出発時刻が経路・順序・後続旅行時間に与える影響を説明できること。\par

\NoteSection{過去問関連}
\ExamItem{WS22/23 Aufgabe 1 (7点)}{TSP の数理最適化モデルと、時間次元を持つTSPで何が変わるかを問う問題。}{標準TSPでは、顧客ノード、アーク、旅行時間または距離、二値変数 x\_ij、各顧客を一度訪問する制約、部分巡回除去制約、総旅行時間最小化を定義する。時間依存TSPでは旅行時間が定数 d\_ij ではなく出発時刻 t に依存する d\_ij(t) となり、あるアーク選択が後続の到着時刻と後続アークの旅行時間を変えるため、静的モデルより決定変数・制約・解法が複雑になる。}
\ExamItem{SS24 Aufgabe 1b (4点)}{旅行時間の情報モデルの例と、常に最も詳細なモデルを選ばない理由を問う問題。}{例として、単一平均旅行時間行列、時間帯別行列、区間旅行時間、確率分布モデルがある。詳細なモデルほど現実に近い可能性はあるが、データ量、信頼性、計算時間、解釈可能性、意思決定モデルの複雑性が悪化する。したがって目的に十分な粒度を選ぶ。}
\ExamItem{WS22/23 Aufgabe 4 (8点)}{過去の交通観測から時間帯別平均旅行時間を構築する手順を説明する問題。}{まずFCDなどから位置・時刻・速度を収集する。位置を道路セグメントにマッチングし、セグメントごとの観測速度を得る。外れ値やエラーを除去し、時間帯 w とセグメント l ごとに観測を集約する。平均速度または旅行時間を計算し、観測不足の部分は類似セグメントや時間帯で一般化する。最後に時間帯別の旅行時間行列を作り、TSP/VRPの情報モデルとして使う。}
\ExamItem{SS24 Aufgabe 1b (4点)}{旅行時間情報モデルの選択と複雑性のトレードオフを問う問題。}{平均値モデルは簡単だがラッシュアワーを無視しやすい。時間帯別モデルは実際の交通変動を反映するが、観測数が少ないセルでは信頼性が低くなる。したがってデータ量と意思決定への効果を見ながら集約・一般化する必要がある。}

\NoteSection{確認問題と解答}
\begin{enumerate}
\item \textbf{時間依存旅行時間モデルでは、なぜ一つの旅行時間行列では不十分か。}\\
一つの行列では全時間帯で同じ旅行時間を仮定することになる。実際には道路セグメントごとに混雑時間が異なり、10時と17時では同じ顧客間の最短経路や所要時間が異なる。そのため時間帯別または連続時刻の情報モデルが必要になる。
\item \textbf{時間依存性がヒューリスティックを難しくする理由を説明せよ。}\\
ヒューリスティックの局所的な選択が、後続の出発時刻を変え、後続の旅行時間関数の評価値を変えるからである。静的な追加距離だけでは全体影響を評価できない。
\end{enumerate}

\end{multicols}

\clearpage

\SlideHeader{Lecture 2: Information and Decision Modeling}{013}{モデリングの概略}

\SlideImage{13}{../Materials/2026/3DM_02_Information-and-Decision-Modeling_SS26.pdf}

\vspace{1.5mm}

\begin{multicols}{2}

\footnotesize

\raggedcolumns

\NoteSection{日本語訳}

\begin{itemize}
\item 情報モデル
\begin{itemize}
\item 顧客間、デポを含む時間依存旅行時間。
\item 時間依存旅行時間行列、または複数の旅行時間行列。
\end{itemize}
\item 意思決定モデル
\begin{itemize}
\item 目的: 旅行時間を最小化する。
\item 制約: 全顧客をサービスする。
\item 制約: ルーティング制約、部分巡回除去など。
\item 決定変数: 時刻 t に顧客 i と j の間を移動するか。
\end{itemize}
\end{itemize}

\NoteSection{解説}
時間依存旅行時間では、顧客 i から j への旅行時間を定数ではなく出発時刻 t の関数として扱う。記号的には d\_ij(t) または tau\_ij(t) と考える。\par

同じ二地点間でも、出発時刻が異なれば最短経路も旅行時間も変わる。したがって、訪問順序の変更は単にアークの足し算を変えるだけでなく、後続の到着時刻と全体の旅行時間を変える。\par

このため、情報モデルは時間帯別行列や時間依存リンク旅行時間を必要とし、意思決定モデルは時間の次元を含む変数・制約を必要とする。\par

\NoteSection{試験対策ポイント}
旅行時間を d\_ij(t) として捉え、出発時刻が経路・順序・後続旅行時間に与える影響を説明できること。\par

\NoteSection{過去問関連}
\ExamItem{WS22/23 Aufgabe 1 (7点)}{TSP の数理最適化モデルと、時間次元を持つTSPで何が変わるかを問う問題。}{標準TSPでは、顧客ノード、アーク、旅行時間または距離、二値変数 x\_ij、各顧客を一度訪問する制約、部分巡回除去制約、総旅行時間最小化を定義する。時間依存TSPでは旅行時間が定数 d\_ij ではなく出発時刻 t に依存する d\_ij(t) となり、あるアーク選択が後続の到着時刻と後続アークの旅行時間を変えるため、静的モデルより決定変数・制約・解法が複雑になる。}
\ExamItem{SS24 Aufgabe 1b (4点)}{旅行時間の情報モデルの例と、常に最も詳細なモデルを選ばない理由を問う問題。}{例として、単一平均旅行時間行列、時間帯別行列、区間旅行時間、確率分布モデルがある。詳細なモデルほど現実に近い可能性はあるが、データ量、信頼性、計算時間、解釈可能性、意思決定モデルの複雑性が悪化する。したがって目的に十分な粒度を選ぶ。}
\ExamItem{WS22/23 Aufgabe 4 (8点)}{過去の交通観測から時間帯別平均旅行時間を構築する手順を説明する問題。}{まずFCDなどから位置・時刻・速度を収集する。位置を道路セグメントにマッチングし、セグメントごとの観測速度を得る。外れ値やエラーを除去し、時間帯 w とセグメント l ごとに観測を集約する。平均速度または旅行時間を計算し、観測不足の部分は類似セグメントや時間帯で一般化する。最後に時間帯別の旅行時間行列を作り、TSP/VRPの情報モデルとして使う。}
\ExamItem{SS24 Aufgabe 1b (4点)}{旅行時間情報モデルの選択と複雑性のトレードオフを問う問題。}{平均値モデルは簡単だがラッシュアワーを無視しやすい。時間帯別モデルは実際の交通変動を反映するが、観測数が少ないセルでは信頼性が低くなる。したがってデータ量と意思決定への効果を見ながら集約・一般化する必要がある。}

\NoteSection{確認問題と解答}
\begin{enumerate}
\item \textbf{時間依存旅行時間モデルでは、なぜ一つの旅行時間行列では不十分か。}\\
一つの行列では全時間帯で同じ旅行時間を仮定することになる。実際には道路セグメントごとに混雑時間が異なり、10時と17時では同じ顧客間の最短経路や所要時間が異なる。そのため時間帯別または連続時刻の情報モデルが必要になる。
\item \textbf{時間依存性がヒューリスティックを難しくする理由を説明せよ。}\\
ヒューリスティックの局所的な選択が、後続の出発時刻を変え、後続の旅行時間関数の評価値を変えるからである。静的な追加距離だけでは全体影響を評価できない。
\end{enumerate}

\end{multicols}

\clearpage

\SlideHeader{Lecture 2: Information and Decision Modeling}{014}{データ収集: Floating Car Data}

\SlideImage{14}{../Materials/2026/3DM_02_Information-and-Decision-Modeling_SS26.pdf}

\vspace{1.5mm}

\begin{multicols}{2}

\footnotesize

\raggedcolumns

\NoteSection{日本語訳}

\begin{itemize}
\item 時刻と道路セグメントごとの旅行時間。
\item 正確な平均値を計算するには大量のデータが必要である。
\item Floating Car Data: 大規模な車両群が、位置、方向、場合によって速度を短い時間間隔で送信する。
\end{itemize}

\NoteSection{解説}
Floating Car Data は、車両群から位置、方向、速度などを高頻度で収集する交通データである。タクシーや商用車のGPSデータが典型である。\par

生データは、そのまま道路リンクごとの速度を与えるわけではない。まず位置を道路セグメントに対応付け、時刻と方向を考慮し、セグメント・時間帯ごとの速度観測へ変換する必要がある。\par

外れ値除去、観測不足の確認、時間帯集約、正規化、クラスタリング、空間・時間検証を通じて、意思決定モデルが使える安定した情報モデルを作る。\par

この一連の処理は、データが先にあり、その後に情報モデルを構築するという Lecture 2 の中心例である。\par

\NoteSection{試験対策ポイント}
FCDの生データから、クリーニング、集約、一般化を経て意思決定モデル用の旅行時間情報を作る手順を説明できること。\par

\NoteSection{過去問関連}
\ExamItem{WS22/23 Aufgabe 4 (8点)}{過去の交通観測から時間帯別平均旅行時間を構築する手順を説明する問題。}{まずFCDなどから位置・時刻・速度を収集する。位置を道路セグメントにマッチングし、セグメントごとの観測速度を得る。外れ値やエラーを除去し、時間帯 w とセグメント l ごとに観測を集約する。平均速度または旅行時間を計算し、観測不足の部分は類似セグメントや時間帯で一般化する。最後に時間帯別の旅行時間行列を作り、TSP/VRPの情報モデルとして使う。}
\ExamItem{SS24 Aufgabe 1b (4点)}{旅行時間情報モデルの選択と複雑性のトレードオフを問う問題。}{平均値モデルは簡単だがラッシュアワーを無視しやすい。時間帯別モデルは実際の交通変動を反映するが、観測数が少ないセルでは信頼性が低くなる。したがってデータ量と意思決定への効果を見ながら集約・一般化する必要がある。}

\NoteSection{確認問題と解答}
\begin{enumerate}
\item \textbf{FCDから時間帯別旅行時間行列を作る手順を説明せよ。}\\
GPS位置と時刻を取得し、道路セグメントにマッチングする。セグメントごとに速度を計算し、外れ値を除去する。曜日・時間帯ごとに観測を集約し、平均速度または旅行時間を計算する。観測不足は一般化・クラスタリングで補い、最後に顧客間の時間依存最短経路を計算して行列化する。
\item \textbf{なぜ正規化とクラスタリングが必要か。}\\
道路ごとに制限速度や絶対速度が違うため、単純な速度値ではパターン比較が難しい。正規化により日内変動の形を比較でき、クラスタリングで似た交通品質パターンを持つセグメントをまとめられる。これにより観測が少ない場所でも一般化しやすい。
\item \textbf{データクリーニングで何を除去すべきか。}\\
速度制限を大きく超える観測、GPSエラー、方向が合わない観測、停止や信号待ちを過度に反映する異常値などを検討する。ただし外部要因もあるため、機械的除去だけでなくモデル目的に合う判断が必要である。
\end{enumerate}

\end{multicols}

\clearpage

\SlideHeader{Lecture 2: Information and Decision Modeling}{015}{FCDの収集と利用}

\SlideImage{15}{../Materials/2026/3DM_02_Information-and-Decision-Modeling_SS26.pdf}

\vspace{1.5mm}

\begin{multicols}{2}

\footnotesize

\raggedcolumns

\NoteSection{日本語訳}

\begin{itemize}
\item タクシー車両群。
\item ディスパッチャー。
\item 公共交通管理。
\item プロバイダ。
\item インターネット。
\item GPRS, SMS, RDS-TMC, HTTP による通信。
\item およそ1分に1回の通信。
\end{itemize}

\NoteSection{解説}
Floating Car Data は、車両群から位置、方向、速度などを高頻度で収集する交通データである。タクシーや商用車のGPSデータが典型である。\par

生データは、そのまま道路リンクごとの速度を与えるわけではない。まず位置を道路セグメントに対応付け、時刻と方向を考慮し、セグメント・時間帯ごとの速度観測へ変換する必要がある。\par

外れ値除去、観測不足の確認、時間帯集約、正規化、クラスタリング、空間・時間検証を通じて、意思決定モデルが使える安定した情報モデルを作る。\par

この一連の処理は、データが先にあり、その後に情報モデルを構築するという Lecture 2 の中心例である。\par

\NoteSection{試験対策ポイント}
FCDの生データから、クリーニング、集約、一般化を経て意思決定モデル用の旅行時間情報を作る手順を説明できること。\par

\NoteSection{過去問関連}
\ExamItem{WS22/23 Aufgabe 4 (8点)}{過去の交通観測から時間帯別平均旅行時間を構築する手順を説明する問題。}{まずFCDなどから位置・時刻・速度を収集する。位置を道路セグメントにマッチングし、セグメントごとの観測速度を得る。外れ値やエラーを除去し、時間帯 w とセグメント l ごとに観測を集約する。平均速度または旅行時間を計算し、観測不足の部分は類似セグメントや時間帯で一般化する。最後に時間帯別の旅行時間行列を作り、TSP/VRPの情報モデルとして使う。}
\ExamItem{SS24 Aufgabe 1b (4点)}{旅行時間情報モデルの選択と複雑性のトレードオフを問う問題。}{平均値モデルは簡単だがラッシュアワーを無視しやすい。時間帯別モデルは実際の交通変動を反映するが、観測数が少ないセルでは信頼性が低くなる。したがってデータ量と意思決定への効果を見ながら集約・一般化する必要がある。}

\NoteSection{確認問題と解答}
\begin{enumerate}
\item \textbf{FCDから時間帯別旅行時間行列を作る手順を説明せよ。}\\
GPS位置と時刻を取得し、道路セグメントにマッチングする。セグメントごとに速度を計算し、外れ値を除去する。曜日・時間帯ごとに観測を集約し、平均速度または旅行時間を計算する。観測不足は一般化・クラスタリングで補い、最後に顧客間の時間依存最短経路を計算して行列化する。
\item \textbf{なぜ正規化とクラスタリングが必要か。}\\
道路ごとに制限速度や絶対速度が違うため、単純な速度値ではパターン比較が難しい。正規化により日内変動の形を比較でき、クラスタリングで似た交通品質パターンを持つセグメントをまとめられる。これにより観測が少ない場所でも一般化しやすい。
\item \textbf{データクリーニングで何を除去すべきか。}\\
速度制限を大きく超える観測、GPSエラー、方向が合わない観測、停止や信号待ちを過度に反映する異常値などを検討する。ただし外部要因もあるため、機械的除去だけでなくモデル目的に合う判断が必要である。
\end{enumerate}

\end{multicols}

\clearpage

\SlideHeader{Lecture 2: Information and Decision Modeling}{016}{FCD: 速度計算}

\SlideImage{16}{../Materials/2026/3DM_02_Information-and-Decision-Modeling_SS26.pdf}

\vspace{1.5mm}

\begin{multicols}{2}

\footnotesize

\raggedcolumns

\NoteSection{日本語訳}

\begin{itemize}
\item タクシーデータは、道路、方向、速度を直接与えない。
\item 必要な手順
\begin{itemize}
\item タクシー位置を観測する。
\item 位置を道路セグメントに投影する。
\item このセグメントの観測速度を計算する。
\end{itemize}
\end{itemize}

\NoteSection{解説}
Floating Car Data は、車両群から位置、方向、速度などを高頻度で収集する交通データである。タクシーや商用車のGPSデータが典型である。\par

生データは、そのまま道路リンクごとの速度を与えるわけではない。まず位置を道路セグメントに対応付け、時刻と方向を考慮し、セグメント・時間帯ごとの速度観測へ変換する必要がある。\par

外れ値除去、観測不足の確認、時間帯集約、正規化、クラスタリング、空間・時間検証を通じて、意思決定モデルが使える安定した情報モデルを作る。\par

この一連の処理は、データが先にあり、その後に情報モデルを構築するという Lecture 2 の中心例である。\par

\NoteSection{試験対策ポイント}
FCDの生データから、クリーニング、集約、一般化を経て意思決定モデル用の旅行時間情報を作る手順を説明できること。\par

\NoteSection{過去問関連}
\ExamItem{WS22/23 Aufgabe 4 (8点)}{過去の交通観測から時間帯別平均旅行時間を構築する手順を説明する問題。}{まずFCDなどから位置・時刻・速度を収集する。位置を道路セグメントにマッチングし、セグメントごとの観測速度を得る。外れ値やエラーを除去し、時間帯 w とセグメント l ごとに観測を集約する。平均速度または旅行時間を計算し、観測不足の部分は類似セグメントや時間帯で一般化する。最後に時間帯別の旅行時間行列を作り、TSP/VRPの情報モデルとして使う。}
\ExamItem{SS24 Aufgabe 1b (4点)}{旅行時間情報モデルの選択と複雑性のトレードオフを問う問題。}{平均値モデルは簡単だがラッシュアワーを無視しやすい。時間帯別モデルは実際の交通変動を反映するが、観測数が少ないセルでは信頼性が低くなる。したがってデータ量と意思決定への効果を見ながら集約・一般化する必要がある。}

\NoteSection{確認問題と解答}
\begin{enumerate}
\item \textbf{FCDから時間帯別旅行時間行列を作る手順を説明せよ。}\\
GPS位置と時刻を取得し、道路セグメントにマッチングする。セグメントごとに速度を計算し、外れ値を除去する。曜日・時間帯ごとに観測を集約し、平均速度または旅行時間を計算する。観測不足は一般化・クラスタリングで補い、最後に顧客間の時間依存最短経路を計算して行列化する。
\item \textbf{なぜ正規化とクラスタリングが必要か。}\\
道路ごとに制限速度や絶対速度が違うため、単純な速度値ではパターン比較が難しい。正規化により日内変動の形を比較でき、クラスタリングで似た交通品質パターンを持つセグメントをまとめられる。これにより観測が少ない場所でも一般化しやすい。
\item \textbf{データクリーニングで何を除去すべきか。}\\
速度制限を大きく超える観測、GPSエラー、方向が合わない観測、停止や信号待ちを過度に反映する異常値などを検討する。ただし外部要因もあるため、機械的除去だけでなくモデル目的に合う判断が必要である。
\end{enumerate}

\end{multicols}

\clearpage

\SlideHeader{Lecture 2: Information and Decision Modeling}{017}{FCD: 速度計算の表}

\SlideImage{17}{../Materials/2026/3DM_02_Information-and-Decision-Modeling_SS26.pdf}

\vspace{1.5mm}

\begin{multicols}{2}

\footnotesize

\raggedcolumns

\NoteSection{日本語訳}

\begin{itemize}
\item link, time, speed の観測表。
\item 同じリンクでも時刻により速度が大きく異なる。
\end{itemize}

\NoteSection{解説}
Floating Car Data は、車両群から位置、方向、速度などを高頻度で収集する交通データである。タクシーや商用車のGPSデータが典型である。\par

生データは、そのまま道路リンクごとの速度を与えるわけではない。まず位置を道路セグメントに対応付け、時刻と方向を考慮し、セグメント・時間帯ごとの速度観測へ変換する必要がある。\par

外れ値除去、観測不足の確認、時間帯集約、正規化、クラスタリング、空間・時間検証を通じて、意思決定モデルが使える安定した情報モデルを作る。\par

この一連の処理は、データが先にあり、その後に情報モデルを構築するという Lecture 2 の中心例である。\par

\NoteSection{試験対策ポイント}
FCDの生データから、クリーニング、集約、一般化を経て意思決定モデル用の旅行時間情報を作る手順を説明できること。\par

\NoteSection{過去問関連}
\ExamItem{WS22/23 Aufgabe 4 (8点)}{過去の交通観測から時間帯別平均旅行時間を構築する手順を説明する問題。}{まずFCDなどから位置・時刻・速度を収集する。位置を道路セグメントにマッチングし、セグメントごとの観測速度を得る。外れ値やエラーを除去し、時間帯 w とセグメント l ごとに観測を集約する。平均速度または旅行時間を計算し、観測不足の部分は類似セグメントや時間帯で一般化する。最後に時間帯別の旅行時間行列を作り、TSP/VRPの情報モデルとして使う。}
\ExamItem{SS24 Aufgabe 1b (4点)}{旅行時間情報モデルの選択と複雑性のトレードオフを問う問題。}{平均値モデルは簡単だがラッシュアワーを無視しやすい。時間帯別モデルは実際の交通変動を反映するが、観測数が少ないセルでは信頼性が低くなる。したがってデータ量と意思決定への効果を見ながら集約・一般化する必要がある。}

\NoteSection{確認問題と解答}
\begin{enumerate}
\item \textbf{FCDから時間帯別旅行時間行列を作る手順を説明せよ。}\\
GPS位置と時刻を取得し、道路セグメントにマッチングする。セグメントごとに速度を計算し、外れ値を除去する。曜日・時間帯ごとに観測を集約し、平均速度または旅行時間を計算する。観測不足は一般化・クラスタリングで補い、最後に顧客間の時間依存最短経路を計算して行列化する。
\item \textbf{なぜ正規化とクラスタリングが必要か。}\\
道路ごとに制限速度や絶対速度が違うため、単純な速度値ではパターン比較が難しい。正規化により日内変動の形を比較でき、クラスタリングで似た交通品質パターンを持つセグメントをまとめられる。これにより観測が少ない場所でも一般化しやすい。
\item \textbf{データクリーニングで何を除去すべきか。}\\
速度制限を大きく超える観測、GPSエラー、方向が合わない観測、停止や信号待ちを過度に反映する異常値などを検討する。ただし外部要因もあるため、機械的除去だけでなくモデル目的に合う判断が必要である。
\end{enumerate}

\end{multicols}

\clearpage

\SlideHeader{Lecture 2: Information and Decision Modeling}{018}{データ分析}

\SlideImage{18}{../Materials/2026/3DM_02_Information-and-Decision-Modeling_SS26.pdf}

\vspace{1.5mm}

\begin{multicols}{2}

\footnotesize

\raggedcolumns

\NoteSection{日本語訳}

\begin{itemize}
\item 一日の中で速度パターンが変化する。これは仮説を確認する。
\item 同じ曜日・同じ時刻でもばらつきがある。これは今回の仮説では無視し、次回扱う。
\item 外れ値やエラーが存在し得る。タクシーデータでは特に注意が必要である。
\end{itemize}

\NoteSection{解説}
Floating Car Data は、車両群から位置、方向、速度などを高頻度で収集する交通データである。タクシーや商用車のGPSデータが典型である。\par

生データは、そのまま道路リンクごとの速度を与えるわけではない。まず位置を道路セグメントに対応付け、時刻と方向を考慮し、セグメント・時間帯ごとの速度観測へ変換する必要がある。\par

外れ値除去、観測不足の確認、時間帯集約、正規化、クラスタリング、空間・時間検証を通じて、意思決定モデルが使える安定した情報モデルを作る。\par

この一連の処理は、データが先にあり、その後に情報モデルを構築するという Lecture 2 の中心例である。\par

\NoteSection{試験対策ポイント}
FCDの生データから、クリーニング、集約、一般化を経て意思決定モデル用の旅行時間情報を作る手順を説明できること。\par

\NoteSection{過去問関連}
\ExamItem{WS22/23 Aufgabe 4 (8点)}{過去の交通観測から時間帯別平均旅行時間を構築する手順を説明する問題。}{まずFCDなどから位置・時刻・速度を収集する。位置を道路セグメントにマッチングし、セグメントごとの観測速度を得る。外れ値やエラーを除去し、時間帯 w とセグメント l ごとに観測を集約する。平均速度または旅行時間を計算し、観測不足の部分は類似セグメントや時間帯で一般化する。最後に時間帯別の旅行時間行列を作り、TSP/VRPの情報モデルとして使う。}
\ExamItem{SS24 Aufgabe 1b (4点)}{旅行時間情報モデルの選択と複雑性のトレードオフを問う問題。}{平均値モデルは簡単だがラッシュアワーを無視しやすい。時間帯別モデルは実際の交通変動を反映するが、観測数が少ないセルでは信頼性が低くなる。したがってデータ量と意思決定への効果を見ながら集約・一般化する必要がある。}

\NoteSection{確認問題と解答}
\begin{enumerate}
\item \textbf{FCDから時間帯別旅行時間行列を作る手順を説明せよ。}\\
GPS位置と時刻を取得し、道路セグメントにマッチングする。セグメントごとに速度を計算し、外れ値を除去する。曜日・時間帯ごとに観測を集約し、平均速度または旅行時間を計算する。観測不足は一般化・クラスタリングで補い、最後に顧客間の時間依存最短経路を計算して行列化する。
\item \textbf{なぜ正規化とクラスタリングが必要か。}\\
道路ごとに制限速度や絶対速度が違うため、単純な速度値ではパターン比較が難しい。正規化により日内変動の形を比較でき、クラスタリングで似た交通品質パターンを持つセグメントをまとめられる。これにより観測が少ない場所でも一般化しやすい。
\item \textbf{データクリーニングで何を除去すべきか。}\\
速度制限を大きく超える観測、GPSエラー、方向が合わない観測、停止や信号待ちを過度に反映する異常値などを検討する。ただし外部要因もあるため、機械的除去だけでなくモデル目的に合う判断が必要である。
\end{enumerate}

\end{multicols}

\clearpage

\SlideHeader{Lecture 2: Information and Decision Modeling}{019}{データクリーニング}

\SlideImage{19}{../Materials/2026/3DM_02_Information-and-Decision-Modeling_SS26.pdf}

\vspace{1.5mm}

\begin{multicols}{2}

\footnotesize

\raggedcolumns

\NoteSection{日本語訳}

\begin{itemize}
\item どの観測が外れ値またはエラーかを判断する。
\item 交通信号、天候など多くの外部要因があるため、客観的評価は難しい。
\item データクリーニング例: 制限速度の1.5倍を超える速度観測をすべてフィルタする。
\end{itemize}

\NoteSection{解説}
Floating Car Data は、車両群から位置、方向、速度などを高頻度で収集する交通データである。タクシーや商用車のGPSデータが典型である。\par

生データは、そのまま道路リンクごとの速度を与えるわけではない。まず位置を道路セグメントに対応付け、時刻と方向を考慮し、セグメント・時間帯ごとの速度観測へ変換する必要がある。\par

外れ値除去、観測不足の確認、時間帯集約、正規化、クラスタリング、空間・時間検証を通じて、意思決定モデルが使える安定した情報モデルを作る。\par

この一連の処理は、データが先にあり、その後に情報モデルを構築するという Lecture 2 の中心例である。\par

\NoteSection{試験対策ポイント}
FCDの生データから、クリーニング、集約、一般化を経て意思決定モデル用の旅行時間情報を作る手順を説明できること。\par

\NoteSection{過去問関連}
\ExamItem{WS22/23 Aufgabe 4 (8点)}{過去の交通観測から時間帯別平均旅行時間を構築する手順を説明する問題。}{まずFCDなどから位置・時刻・速度を収集する。位置を道路セグメントにマッチングし、セグメントごとの観測速度を得る。外れ値やエラーを除去し、時間帯 w とセグメント l ごとに観測を集約する。平均速度または旅行時間を計算し、観測不足の部分は類似セグメントや時間帯で一般化する。最後に時間帯別の旅行時間行列を作り、TSP/VRPの情報モデルとして使う。}
\ExamItem{SS24 Aufgabe 1b (4点)}{旅行時間情報モデルの選択と複雑性のトレードオフを問う問題。}{平均値モデルは簡単だがラッシュアワーを無視しやすい。時間帯別モデルは実際の交通変動を反映するが、観測数が少ないセルでは信頼性が低くなる。したがってデータ量と意思決定への効果を見ながら集約・一般化する必要がある。}

\NoteSection{確認問題と解答}
\begin{enumerate}
\item \textbf{FCDから時間帯別旅行時間行列を作る手順を説明せよ。}\\
GPS位置と時刻を取得し、道路セグメントにマッチングする。セグメントごとに速度を計算し、外れ値を除去する。曜日・時間帯ごとに観測を集約し、平均速度または旅行時間を計算する。観測不足は一般化・クラスタリングで補い、最後に顧客間の時間依存最短経路を計算して行列化する。
\item \textbf{なぜ正規化とクラスタリングが必要か。}\\
道路ごとに制限速度や絶対速度が違うため、単純な速度値ではパターン比較が難しい。正規化により日内変動の形を比較でき、クラスタリングで似た交通品質パターンを持つセグメントをまとめられる。これにより観測が少ない場所でも一般化しやすい。
\item \textbf{データクリーニングで何を除去すべきか。}\\
速度制限を大きく超える観測、GPSエラー、方向が合わない観測、停止や信号待ちを過度に反映する異常値などを検討する。ただし外部要因もあるため、機械的除去だけでなくモデル目的に合う判断が必要である。
\end{enumerate}

\end{multicols}

\clearpage

\SlideHeader{Lecture 2: Information and Decision Modeling}{020}{データ可視化}

\SlideImage{20}{../Materials/2026/3DM_02_Information-and-Decision-Modeling_SS26.pdf}

\vspace{1.5mm}

\begin{multicols}{2}

\footnotesize

\raggedcolumns

\NoteSection{日本語訳}

\begin{itemize}
\item 観測の空間分布を見る。
\item よくカバーされている地域がある。
\item 観測が少ない地域もある。これはタクシーデータに由来する。
\item 時間帯についても同様である。
\item 観測を信頼できるか。
\item 集約と一般化が必要である。
\end{itemize}

\NoteSection{解説}
Floating Car Data は、車両群から位置、方向、速度などを高頻度で収集する交通データである。タクシーや商用車のGPSデータが典型である。\par

生データは、そのまま道路リンクごとの速度を与えるわけではない。まず位置を道路セグメントに対応付け、時刻と方向を考慮し、セグメント・時間帯ごとの速度観測へ変換する必要がある。\par

外れ値除去、観測不足の確認、時間帯集約、正規化、クラスタリング、空間・時間検証を通じて、意思決定モデルが使える安定した情報モデルを作る。\par

この一連の処理は、データが先にあり、その後に情報モデルを構築するという Lecture 2 の中心例である。\par

\NoteSection{試験対策ポイント}
FCDの生データから、クリーニング、集約、一般化を経て意思決定モデル用の旅行時間情報を作る手順を説明できること。\par

\NoteSection{過去問関連}
\ExamItem{WS22/23 Aufgabe 4 (8点)}{過去の交通観測から時間帯別平均旅行時間を構築する手順を説明する問題。}{まずFCDなどから位置・時刻・速度を収集する。位置を道路セグメントにマッチングし、セグメントごとの観測速度を得る。外れ値やエラーを除去し、時間帯 w とセグメント l ごとに観測を集約する。平均速度または旅行時間を計算し、観測不足の部分は類似セグメントや時間帯で一般化する。最後に時間帯別の旅行時間行列を作り、TSP/VRPの情報モデルとして使う。}
\ExamItem{SS24 Aufgabe 1b (4点)}{旅行時間情報モデルの選択と複雑性のトレードオフを問う問題。}{平均値モデルは簡単だがラッシュアワーを無視しやすい。時間帯別モデルは実際の交通変動を反映するが、観測数が少ないセルでは信頼性が低くなる。したがってデータ量と意思決定への効果を見ながら集約・一般化する必要がある。}

\NoteSection{確認問題と解答}
\begin{enumerate}
\item \textbf{FCDから時間帯別旅行時間行列を作る手順を説明せよ。}\\
GPS位置と時刻を取得し、道路セグメントにマッチングする。セグメントごとに速度を計算し、外れ値を除去する。曜日・時間帯ごとに観測を集約し、平均速度または旅行時間を計算する。観測不足は一般化・クラスタリングで補い、最後に顧客間の時間依存最短経路を計算して行列化する。
\item \textbf{なぜ正規化とクラスタリングが必要か。}\\
道路ごとに制限速度や絶対速度が違うため、単純な速度値ではパターン比較が難しい。正規化により日内変動の形を比較でき、クラスタリングで似た交通品質パターンを持つセグメントをまとめられる。これにより観測が少ない場所でも一般化しやすい。
\item \textbf{データクリーニングで何を除去すべきか。}\\
速度制限を大きく超える観測、GPSエラー、方向が合わない観測、停止や信号待ちを過度に反映する異常値などを検討する。ただし外部要因もあるため、機械的除去だけでなくモデル目的に合う判断が必要である。
\end{enumerate}

\end{multicols}

\clearpage

\SlideHeader{Lecture 2: Information and Decision Modeling}{021}{集約}

\SlideImage{21}{../Materials/2026/3DM_02_Information-and-Decision-Modeling_SS26.pdf}

\vspace{1.5mm}

\begin{multicols}{2}

\footnotesize

\raggedcolumns

\NoteSection{日本語訳}

\begin{itemize}
\item 時間ステップ W を定義する。例: 24×7。
\item セグメント l、時間ステップ w ごとに、観測速度から平均旅行速度を計算する。
\item W=1: セグメントごとに一つの値。
\item W=24×7: 一日あたり24個、週7日分の値。
\end{itemize}

\NoteSection{解説}
Floating Car Data は、車両群から位置、方向、速度などを高頻度で収集する交通データである。タクシーや商用車のGPSデータが典型である。\par

生データは、そのまま道路リンクごとの速度を与えるわけではない。まず位置を道路セグメントに対応付け、時刻と方向を考慮し、セグメント・時間帯ごとの速度観測へ変換する必要がある。\par

外れ値除去、観測不足の確認、時間帯集約、正規化、クラスタリング、空間・時間検証を通じて、意思決定モデルが使える安定した情報モデルを作る。\par

この一連の処理は、データが先にあり、その後に情報モデルを構築するという Lecture 2 の中心例である。\par

\NoteSection{試験対策ポイント}
FCDの生データから、クリーニング、集約、一般化を経て意思決定モデル用の旅行時間情報を作る手順を説明できること。\par

\NoteSection{過去問関連}
\ExamItem{WS22/23 Aufgabe 4 (8点)}{過去の交通観測から時間帯別平均旅行時間を構築する手順を説明する問題。}{まずFCDなどから位置・時刻・速度を収集する。位置を道路セグメントにマッチングし、セグメントごとの観測速度を得る。外れ値やエラーを除去し、時間帯 w とセグメント l ごとに観測を集約する。平均速度または旅行時間を計算し、観測不足の部分は類似セグメントや時間帯で一般化する。最後に時間帯別の旅行時間行列を作り、TSP/VRPの情報モデルとして使う。}
\ExamItem{SS24 Aufgabe 1b (4点)}{旅行時間情報モデルの選択と複雑性のトレードオフを問う問題。}{平均値モデルは簡単だがラッシュアワーを無視しやすい。時間帯別モデルは実際の交通変動を反映するが、観測数が少ないセルでは信頼性が低くなる。したがってデータ量と意思決定への効果を見ながら集約・一般化する必要がある。}

\NoteSection{確認問題と解答}
\begin{enumerate}
\item \textbf{FCDから時間帯別旅行時間行列を作る手順を説明せよ。}\\
GPS位置と時刻を取得し、道路セグメントにマッチングする。セグメントごとに速度を計算し、外れ値を除去する。曜日・時間帯ごとに観測を集約し、平均速度または旅行時間を計算する。観測不足は一般化・クラスタリングで補い、最後に顧客間の時間依存最短経路を計算して行列化する。
\item \textbf{なぜ正規化とクラスタリングが必要か。}\\
道路ごとに制限速度や絶対速度が違うため、単純な速度値ではパターン比較が難しい。正規化により日内変動の形を比較でき、クラスタリングで似た交通品質パターンを持つセグメントをまとめられる。これにより観測が少ない場所でも一般化しやすい。
\item \textbf{データクリーニングで何を除去すべきか。}\\
速度制限を大きく超える観測、GPSエラー、方向が合わない観測、停止や信号待ちを過度に反映する異常値などを検討する。ただし外部要因もあるため、機械的除去だけでなくモデル目的に合う判断が必要である。
\end{enumerate}

\end{multicols}

\clearpage

\SlideHeader{Lecture 2: Information and Decision Modeling}{022}{一般化}

\SlideImage{22}{../Materials/2026/3DM_02_Information-and-Decision-Modeling_SS26.pdf}

\vspace{1.5mm}

\begin{multicols}{2}

\footnotesize

\raggedcolumns

\NoteSection{日本語訳}

\begin{itemize}
\item 集約結果は数百万個の値になる。
\item 一部の値は少数の観測にしか基づかない。
\item 意思決定モデルは、アクセスしやすく信頼できる情報モデルを必要とする。
\item データ可視化は、異なる道路セグメントで類似した交通品質パターンを示す。
\item 考え方: 典型的な交通品質に基づいて道路セグメントをグループ化する。
\end{itemize}

\NoteSection{解説}
Floating Car Data は、車両群から位置、方向、速度などを高頻度で収集する交通データである。タクシーや商用車のGPSデータが典型である。\par

生データは、そのまま道路リンクごとの速度を与えるわけではない。まず位置を道路セグメントに対応付け、時刻と方向を考慮し、セグメント・時間帯ごとの速度観測へ変換する必要がある。\par

外れ値除去、観測不足の確認、時間帯集約、正規化、クラスタリング、空間・時間検証を通じて、意思決定モデルが使える安定した情報モデルを作る。\par

この一連の処理は、データが先にあり、その後に情報モデルを構築するという Lecture 2 の中心例である。\par

\NoteSection{試験対策ポイント}
FCDの生データから、クリーニング、集約、一般化を経て意思決定モデル用の旅行時間情報を作る手順を説明できること。\par

\NoteSection{過去問関連}
\ExamItem{WS22/23 Aufgabe 4 (8点)}{過去の交通観測から時間帯別平均旅行時間を構築する手順を説明する問題。}{まずFCDなどから位置・時刻・速度を収集する。位置を道路セグメントにマッチングし、セグメントごとの観測速度を得る。外れ値やエラーを除去し、時間帯 w とセグメント l ごとに観測を集約する。平均速度または旅行時間を計算し、観測不足の部分は類似セグメントや時間帯で一般化する。最後に時間帯別の旅行時間行列を作り、TSP/VRPの情報モデルとして使う。}
\ExamItem{SS24 Aufgabe 1b (4点)}{旅行時間情報モデルの選択と複雑性のトレードオフを問う問題。}{平均値モデルは簡単だがラッシュアワーを無視しやすい。時間帯別モデルは実際の交通変動を反映するが、観測数が少ないセルでは信頼性が低くなる。したがってデータ量と意思決定への効果を見ながら集約・一般化する必要がある。}

\NoteSection{確認問題と解答}
\begin{enumerate}
\item \textbf{FCDから時間帯別旅行時間行列を作る手順を説明せよ。}\\
GPS位置と時刻を取得し、道路セグメントにマッチングする。セグメントごとに速度を計算し、外れ値を除去する。曜日・時間帯ごとに観測を集約し、平均速度または旅行時間を計算する。観測不足は一般化・クラスタリングで補い、最後に顧客間の時間依存最短経路を計算して行列化する。
\item \textbf{なぜ正規化とクラスタリングが必要か。}\\
道路ごとに制限速度や絶対速度が違うため、単純な速度値ではパターン比較が難しい。正規化により日内変動の形を比較でき、クラスタリングで似た交通品質パターンを持つセグメントをまとめられる。これにより観測が少ない場所でも一般化しやすい。
\item \textbf{データクリーニングで何を除去すべきか。}\\
速度制限を大きく超える観測、GPSエラー、方向が合わない観測、停止や信号待ちを過度に反映する異常値などを検討する。ただし外部要因もあるため、機械的除去だけでなくモデル目的に合う判断が必要である。
\end{enumerate}

\end{multicols}

\clearpage

\SlideHeader{Lecture 2: Information and Decision Modeling}{023}{一般化: 正規化}

\SlideImage{23}{../Materials/2026/3DM_02_Information-and-Decision-Modeling_SS26.pdf}

\vspace{1.5mm}

\begin{multicols}{2}

\footnotesize

\raggedcolumns

\NoteSection{日本語訳}

\begin{itemize}
\item 異なるセグメントは同じパターンを持っていても絶対値が異なる。例: 制限速度30または50。
\item 比較を可能にするため、日内最大値で正規化する。
\end{itemize}

\NoteSection{解説}
Floating Car Data は、車両群から位置、方向、速度などを高頻度で収集する交通データである。タクシーや商用車のGPSデータが典型である。\par

生データは、そのまま道路リンクごとの速度を与えるわけではない。まず位置を道路セグメントに対応付け、時刻と方向を考慮し、セグメント・時間帯ごとの速度観測へ変換する必要がある。\par

外れ値除去、観測不足の確認、時間帯集約、正規化、クラスタリング、空間・時間検証を通じて、意思決定モデルが使える安定した情報モデルを作る。\par

この一連の処理は、データが先にあり、その後に情報モデルを構築するという Lecture 2 の中心例である。\par

\NoteSection{試験対策ポイント}
FCDの生データから、クリーニング、集約、一般化を経て意思決定モデル用の旅行時間情報を作る手順を説明できること。\par

\NoteSection{過去問関連}
\ExamItem{WS22/23 Aufgabe 4 (8点)}{過去の交通観測から時間帯別平均旅行時間を構築する手順を説明する問題。}{まずFCDなどから位置・時刻・速度を収集する。位置を道路セグメントにマッチングし、セグメントごとの観測速度を得る。外れ値やエラーを除去し、時間帯 w とセグメント l ごとに観測を集約する。平均速度または旅行時間を計算し、観測不足の部分は類似セグメントや時間帯で一般化する。最後に時間帯別の旅行時間行列を作り、TSP/VRPの情報モデルとして使う。}
\ExamItem{SS24 Aufgabe 1b (4点)}{旅行時間情報モデルの選択と複雑性のトレードオフを問う問題。}{平均値モデルは簡単だがラッシュアワーを無視しやすい。時間帯別モデルは実際の交通変動を反映するが、観測数が少ないセルでは信頼性が低くなる。したがってデータ量と意思決定への効果を見ながら集約・一般化する必要がある。}

\NoteSection{確認問題と解答}
\begin{enumerate}
\item \textbf{FCDから時間帯別旅行時間行列を作る手順を説明せよ。}\\
GPS位置と時刻を取得し、道路セグメントにマッチングする。セグメントごとに速度を計算し、外れ値を除去する。曜日・時間帯ごとに観測を集約し、平均速度または旅行時間を計算する。観測不足は一般化・クラスタリングで補い、最後に顧客間の時間依存最短経路を計算して行列化する。
\item \textbf{なぜ正規化とクラスタリングが必要か。}\\
道路ごとに制限速度や絶対速度が違うため、単純な速度値ではパターン比較が難しい。正規化により日内変動の形を比較でき、クラスタリングで似た交通品質パターンを持つセグメントをまとめられる。これにより観測が少ない場所でも一般化しやすい。
\item \textbf{データクリーニングで何を除去すべきか。}\\
速度制限を大きく超える観測、GPSエラー、方向が合わない観測、停止や信号待ちを過度に反映する異常値などを検討する。ただし外部要因もあるため、機械的除去だけでなくモデル目的に合う判断が必要である。
\end{enumerate}

\end{multicols}

\clearpage

\SlideHeader{Lecture 2: Information and Decision Modeling}{024}{一般化: クラスタリング}

\SlideImage{24}{../Materials/2026/3DM_02_Information-and-Decision-Modeling_SS26.pdf}

\vspace{1.5mm}

\begin{multicols}{2}

\footnotesize

\raggedcolumns

\NoteSection{日本語訳}

\begin{itemize}
\item クラスタリングアルゴリズムで類似した日内変化をグループ化する。
\item 6つのクラスタ。
\item 0は交通品質が悪く旅行時間が長い。
\item 1は交通品質が良く旅行時間が短い。
\end{itemize}

\NoteSection{解説}
Floating Car Data は、車両群から位置、方向、速度などを高頻度で収集する交通データである。タクシーや商用車のGPSデータが典型である。\par

生データは、そのまま道路リンクごとの速度を与えるわけではない。まず位置を道路セグメントに対応付け、時刻と方向を考慮し、セグメント・時間帯ごとの速度観測へ変換する必要がある。\par

外れ値除去、観測不足の確認、時間帯集約、正規化、クラスタリング、空間・時間検証を通じて、意思決定モデルが使える安定した情報モデルを作る。\par

この一連の処理は、データが先にあり、その後に情報モデルを構築するという Lecture 2 の中心例である。\par

\NoteSection{試験対策ポイント}
FCDの生データから、クリーニング、集約、一般化を経て意思決定モデル用の旅行時間情報を作る手順を説明できること。\par

\NoteSection{過去問関連}
\ExamItem{WS22/23 Aufgabe 4 (8点)}{過去の交通観測から時間帯別平均旅行時間を構築する手順を説明する問題。}{まずFCDなどから位置・時刻・速度を収集する。位置を道路セグメントにマッチングし、セグメントごとの観測速度を得る。外れ値やエラーを除去し、時間帯 w とセグメント l ごとに観測を集約する。平均速度または旅行時間を計算し、観測不足の部分は類似セグメントや時間帯で一般化する。最後に時間帯別の旅行時間行列を作り、TSP/VRPの情報モデルとして使う。}
\ExamItem{SS24 Aufgabe 1b (4点)}{旅行時間情報モデルの選択と複雑性のトレードオフを問う問題。}{平均値モデルは簡単だがラッシュアワーを無視しやすい。時間帯別モデルは実際の交通変動を反映するが、観測数が少ないセルでは信頼性が低くなる。したがってデータ量と意思決定への効果を見ながら集約・一般化する必要がある。}

\NoteSection{確認問題と解答}
\begin{enumerate}
\item \textbf{FCDから時間帯別旅行時間行列を作る手順を説明せよ。}\\
GPS位置と時刻を取得し、道路セグメントにマッチングする。セグメントごとに速度を計算し、外れ値を除去する。曜日・時間帯ごとに観測を集約し、平均速度または旅行時間を計算する。観測不足は一般化・クラスタリングで補い、最後に顧客間の時間依存最短経路を計算して行列化する。
\item \textbf{なぜ正規化とクラスタリングが必要か。}\\
道路ごとに制限速度や絶対速度が違うため、単純な速度値ではパターン比較が難しい。正規化により日内変動の形を比較でき、クラスタリングで似た交通品質パターンを持つセグメントをまとめられる。これにより観測が少ない場所でも一般化しやすい。
\item \textbf{データクリーニングで何を除去すべきか。}\\
速度制限を大きく超える観測、GPSエラー、方向が合わない観測、停止や信号待ちを過度に反映する異常値などを検討する。ただし外部要因もあるため、機械的除去だけでなくモデル目的に合う判断が必要である。
\end{enumerate}

\end{multicols}

\clearpage

\SlideHeader{Lecture 2: Information and Decision Modeling}{025}{空間的検証}

\SlideImage{25}{../Materials/2026/3DM_02_Information-and-Decision-Modeling_SS26.pdf}

\vspace{1.5mm}

\begin{multicols}{2}

\footnotesize

\raggedcolumns

\NoteSection{日本語訳}

\begin{itemize}
\item クラスタはピーク時の挙動が異なる。
\item クラスタは異なる道路セグメントタイプを表す。
\begin{itemize}
\item 市中心部
\item 商業地区
\item 高速道路
\item コンベヤーベルト型の道路
\end{itemize}
\end{itemize}

\NoteSection{解説}
Floating Car Data は、車両群から位置、方向、速度などを高頻度で収集する交通データである。タクシーや商用車のGPSデータが典型である。\par

生データは、そのまま道路リンクごとの速度を与えるわけではない。まず位置を道路セグメントに対応付け、時刻と方向を考慮し、セグメント・時間帯ごとの速度観測へ変換する必要がある。\par

外れ値除去、観測不足の確認、時間帯集約、正規化、クラスタリング、空間・時間検証を通じて、意思決定モデルが使える安定した情報モデルを作る。\par

この一連の処理は、データが先にあり、その後に情報モデルを構築するという Lecture 2 の中心例である。\par

\NoteSection{試験対策ポイント}
FCDの生データから、クリーニング、集約、一般化を経て意思決定モデル用の旅行時間情報を作る手順を説明できること。\par

\NoteSection{過去問関連}
\ExamItem{WS22/23 Aufgabe 4 (8点)}{過去の交通観測から時間帯別平均旅行時間を構築する手順を説明する問題。}{まずFCDなどから位置・時刻・速度を収集する。位置を道路セグメントにマッチングし、セグメントごとの観測速度を得る。外れ値やエラーを除去し、時間帯 w とセグメント l ごとに観測を集約する。平均速度または旅行時間を計算し、観測不足の部分は類似セグメントや時間帯で一般化する。最後に時間帯別の旅行時間行列を作り、TSP/VRPの情報モデルとして使う。}
\ExamItem{SS24 Aufgabe 1b (4点)}{旅行時間情報モデルの選択と複雑性のトレードオフを問う問題。}{平均値モデルは簡単だがラッシュアワーを無視しやすい。時間帯別モデルは実際の交通変動を反映するが、観測数が少ないセルでは信頼性が低くなる。したがってデータ量と意思決定への効果を見ながら集約・一般化する必要がある。}

\NoteSection{確認問題と解答}
\begin{enumerate}
\item \textbf{FCDから時間帯別旅行時間行列を作る手順を説明せよ。}\\
GPS位置と時刻を取得し、道路セグメントにマッチングする。セグメントごとに速度を計算し、外れ値を除去する。曜日・時間帯ごとに観測を集約し、平均速度または旅行時間を計算する。観測不足は一般化・クラスタリングで補い、最後に顧客間の時間依存最短経路を計算して行列化する。
\item \textbf{なぜ正規化とクラスタリングが必要か。}\\
道路ごとに制限速度や絶対速度が違うため、単純な速度値ではパターン比較が難しい。正規化により日内変動の形を比較でき、クラスタリングで似た交通品質パターンを持つセグメントをまとめられる。これにより観測が少ない場所でも一般化しやすい。
\item \textbf{データクリーニングで何を除去すべきか。}\\
速度制限を大きく超える観測、GPSエラー、方向が合わない観測、停止や信号待ちを過度に反映する異常値などを検討する。ただし外部要因もあるため、機械的除去だけでなくモデル目的に合う判断が必要である。
\end{enumerate}

\end{multicols}

\clearpage

\SlideHeader{Lecture 2: Information and Decision Modeling}{026}{時間的検証}

\SlideImage{26}{../Materials/2026/3DM_02_Information-and-Decision-Modeling_SS26.pdf}

\vspace{1.5mm}

\begin{multicols}{2}

\footnotesize

\raggedcolumns

\NoteSection{日本語訳}

\begin{itemize}
\item 一日の中の交通品質。Aは高品質、Fは低品質。
\end{itemize}

\NoteSection{解説}
Floating Car Data は、車両群から位置、方向、速度などを高頻度で収集する交通データである。タクシーや商用車のGPSデータが典型である。\par

生データは、そのまま道路リンクごとの速度を与えるわけではない。まず位置を道路セグメントに対応付け、時刻と方向を考慮し、セグメント・時間帯ごとの速度観測へ変換する必要がある。\par

外れ値除去、観測不足の確認、時間帯集約、正規化、クラスタリング、空間・時間検証を通じて、意思決定モデルが使える安定した情報モデルを作る。\par

この一連の処理は、データが先にあり、その後に情報モデルを構築するという Lecture 2 の中心例である。\par

\NoteSection{試験対策ポイント}
FCDの生データから、クリーニング、集約、一般化を経て意思決定モデル用の旅行時間情報を作る手順を説明できること。\par

\NoteSection{過去問関連}
\ExamItem{WS22/23 Aufgabe 4 (8点)}{過去の交通観測から時間帯別平均旅行時間を構築する手順を説明する問題。}{まずFCDなどから位置・時刻・速度を収集する。位置を道路セグメントにマッチングし、セグメントごとの観測速度を得る。外れ値やエラーを除去し、時間帯 w とセグメント l ごとに観測を集約する。平均速度または旅行時間を計算し、観測不足の部分は類似セグメントや時間帯で一般化する。最後に時間帯別の旅行時間行列を作り、TSP/VRPの情報モデルとして使う。}
\ExamItem{SS24 Aufgabe 1b (4点)}{旅行時間情報モデルの選択と複雑性のトレードオフを問う問題。}{平均値モデルは簡単だがラッシュアワーを無視しやすい。時間帯別モデルは実際の交通変動を反映するが、観測数が少ないセルでは信頼性が低くなる。したがってデータ量と意思決定への効果を見ながら集約・一般化する必要がある。}

\NoteSection{確認問題と解答}
\begin{enumerate}
\item \textbf{FCDから時間帯別旅行時間行列を作る手順を説明せよ。}\\
GPS位置と時刻を取得し、道路セグメントにマッチングする。セグメントごとに速度を計算し、外れ値を除去する。曜日・時間帯ごとに観測を集約し、平均速度または旅行時間を計算する。観測不足は一般化・クラスタリングで補い、最後に顧客間の時間依存最短経路を計算して行列化する。
\item \textbf{なぜ正規化とクラスタリングが必要か。}\\
道路ごとに制限速度や絶対速度が違うため、単純な速度値ではパターン比較が難しい。正規化により日内変動の形を比較でき、クラスタリングで似た交通品質パターンを持つセグメントをまとめられる。これにより観測が少ない場所でも一般化しやすい。
\item \textbf{データクリーニングで何を除去すべきか。}\\
速度制限を大きく超える観測、GPSエラー、方向が合わない観測、停止や信号待ちを過度に反映する異常値などを検討する。ただし外部要因もあるため、機械的除去だけでなくモデル目的に合う判断が必要である。
\end{enumerate}

\end{multicols}

\clearpage

\SlideHeader{Lecture 2: Information and Decision Modeling}{027}{空間・時間: 深夜}

\SlideImage{27}{../Materials/2026/3DM_02_Information-and-Decision-Modeling_SS26.pdf}

\vspace{1.5mm}

\begin{multicols}{2}

\footnotesize

\raggedcolumns

\NoteSection{日本語訳}

\begin{itemize}
\item 月曜日3-4時。
\item セグメントの90\%は空いている。
\item 自由流交通。
\end{itemize}

\NoteSection{解説}
Floating Car Data は、車両群から位置、方向、速度などを高頻度で収集する交通データである。タクシーや商用車のGPSデータが典型である。\par

生データは、そのまま道路リンクごとの速度を与えるわけではない。まず位置を道路セグメントに対応付け、時刻と方向を考慮し、セグメント・時間帯ごとの速度観測へ変換する必要がある。\par

外れ値除去、観測不足の確認、時間帯集約、正規化、クラスタリング、空間・時間検証を通じて、意思決定モデルが使える安定した情報モデルを作る。\par

この一連の処理は、データが先にあり、その後に情報モデルを構築するという Lecture 2 の中心例である。\par

\NoteSection{試験対策ポイント}
FCDの生データから、クリーニング、集約、一般化を経て意思決定モデル用の旅行時間情報を作る手順を説明できること。\par

\NoteSection{過去問関連}
\ExamItem{WS22/23 Aufgabe 4 (8点)}{過去の交通観測から時間帯別平均旅行時間を構築する手順を説明する問題。}{まずFCDなどから位置・時刻・速度を収集する。位置を道路セグメントにマッチングし、セグメントごとの観測速度を得る。外れ値やエラーを除去し、時間帯 w とセグメント l ごとに観測を集約する。平均速度または旅行時間を計算し、観測不足の部分は類似セグメントや時間帯で一般化する。最後に時間帯別の旅行時間行列を作り、TSP/VRPの情報モデルとして使う。}
\ExamItem{SS24 Aufgabe 1b (4点)}{旅行時間情報モデルの選択と複雑性のトレードオフを問う問題。}{平均値モデルは簡単だがラッシュアワーを無視しやすい。時間帯別モデルは実際の交通変動を反映するが、観測数が少ないセルでは信頼性が低くなる。したがってデータ量と意思決定への効果を見ながら集約・一般化する必要がある。}

\NoteSection{確認問題と解答}
\begin{enumerate}
\item \textbf{FCDから時間帯別旅行時間行列を作る手順を説明せよ。}\\
GPS位置と時刻を取得し、道路セグメントにマッチングする。セグメントごとに速度を計算し、外れ値を除去する。曜日・時間帯ごとに観測を集約し、平均速度または旅行時間を計算する。観測不足は一般化・クラスタリングで補い、最後に顧客間の時間依存最短経路を計算して行列化する。
\item \textbf{なぜ正規化とクラスタリングが必要か。}\\
道路ごとに制限速度や絶対速度が違うため、単純な速度値ではパターン比較が難しい。正規化により日内変動の形を比較でき、クラスタリングで似た交通品質パターンを持つセグメントをまとめられる。これにより観測が少ない場所でも一般化しやすい。
\item \textbf{データクリーニングで何を除去すべきか。}\\
速度制限を大きく超える観測、GPSエラー、方向が合わない観測、停止や信号待ちを過度に反映する異常値などを検討する。ただし外部要因もあるため、機械的除去だけでなくモデル目的に合う判断が必要である。
\end{enumerate}

\end{multicols}

\clearpage

\SlideHeader{Lecture 2: Information and Decision Modeling}{028}{空間・時間: ラッシュアワー}

\SlideImage{28}{../Materials/2026/3DM_02_Information-and-Decision-Modeling_SS26.pdf}

\vspace{1.5mm}

\begin{multicols}{2}

\footnotesize

\raggedcolumns

\NoteSection{日本語訳}

\begin{itemize}
\item 月曜日17-18時。
\item 30\%が混雑している。
\item ラッシュアワー。
\end{itemize}

\NoteSection{解説}
Floating Car Data は、車両群から位置、方向、速度などを高頻度で収集する交通データである。タクシーや商用車のGPSデータが典型である。\par

生データは、そのまま道路リンクごとの速度を与えるわけではない。まず位置を道路セグメントに対応付け、時刻と方向を考慮し、セグメント・時間帯ごとの速度観測へ変換する必要がある。\par

外れ値除去、観測不足の確認、時間帯集約、正規化、クラスタリング、空間・時間検証を通じて、意思決定モデルが使える安定した情報モデルを作る。\par

この一連の処理は、データが先にあり、その後に情報モデルを構築するという Lecture 2 の中心例である。\par

\NoteSection{試験対策ポイント}
FCDの生データから、クリーニング、集約、一般化を経て意思決定モデル用の旅行時間情報を作る手順を説明できること。\par

\NoteSection{過去問関連}
\ExamItem{WS22/23 Aufgabe 4 (8点)}{過去の交通観測から時間帯別平均旅行時間を構築する手順を説明する問題。}{まずFCDなどから位置・時刻・速度を収集する。位置を道路セグメントにマッチングし、セグメントごとの観測速度を得る。外れ値やエラーを除去し、時間帯 w とセグメント l ごとに観測を集約する。平均速度または旅行時間を計算し、観測不足の部分は類似セグメントや時間帯で一般化する。最後に時間帯別の旅行時間行列を作り、TSP/VRPの情報モデルとして使う。}
\ExamItem{SS24 Aufgabe 1b (4点)}{旅行時間情報モデルの選択と複雑性のトレードオフを問う問題。}{平均値モデルは簡単だがラッシュアワーを無視しやすい。時間帯別モデルは実際の交通変動を反映するが、観測数が少ないセルでは信頼性が低くなる。したがってデータ量と意思決定への効果を見ながら集約・一般化する必要がある。}

\NoteSection{確認問題と解答}
\begin{enumerate}
\item \textbf{FCDから時間帯別旅行時間行列を作る手順を説明せよ。}\\
GPS位置と時刻を取得し、道路セグメントにマッチングする。セグメントごとに速度を計算し、外れ値を除去する。曜日・時間帯ごとに観測を集約し、平均速度または旅行時間を計算する。観測不足は一般化・クラスタリングで補い、最後に顧客間の時間依存最短経路を計算して行列化する。
\item \textbf{なぜ正規化とクラスタリングが必要か。}\\
道路ごとに制限速度や絶対速度が違うため、単純な速度値ではパターン比較が難しい。正規化により日内変動の形を比較でき、クラスタリングで似た交通品質パターンを持つセグメントをまとめられる。これにより観測が少ない場所でも一般化しやすい。
\item \textbf{データクリーニングで何を除去すべきか。}\\
速度制限を大きく超える観測、GPSエラー、方向が合わない観測、停止や信号待ちを過度に反映する異常値などを検討する。ただし外部要因もあるため、機械的除去だけでなくモデル目的に合う判断が必要である。
\end{enumerate}

\end{multicols}

\clearpage

\SlideHeader{Lecture 2: Information and Decision Modeling}{029}{情報モデル: 結論}

\SlideImage{29}{../Materials/2026/3DM_02_Information-and-Decision-Modeling_SS26.pdf}

\vspace{1.5mm}

\begin{multicols}{2}

\footnotesize

\raggedcolumns

\NoteSection{日本語訳}

\begin{itemize}
\item 道路セグメントと各時間について旅行時間行列を持つ。
\item 例: 最初の時間は長い旅行時間、二番目の時間は短い旅行時間、三番目の時間は最長の旅行時間。
\item この情報に基づいてどのようにルーティングするか。
\end{itemize}

\NoteSection{解説}
時間依存旅行時間では、顧客 i から j への旅行時間を定数ではなく出発時刻 t の関数として扱う。記号的には d\_ij(t) または tau\_ij(t) と考える。\par

同じ二地点間でも、出発時刻が異なれば最短経路も旅行時間も変わる。したがって、訪問順序の変更は単にアークの足し算を変えるだけでなく、後続の到着時刻と全体の旅行時間を変える。\par

このため、情報モデルは時間帯別行列や時間依存リンク旅行時間を必要とし、意思決定モデルは時間の次元を含む変数・制約を必要とする。\par

\NoteSection{試験対策ポイント}
旅行時間を d\_ij(t) として捉え、出発時刻が経路・順序・後続旅行時間に与える影響を説明できること。\par

\NoteSection{過去問関連}
\ExamItem{WS22/23 Aufgabe 1 (7点)}{TSP の数理最適化モデルと、時間次元を持つTSPで何が変わるかを問う問題。}{標準TSPでは、顧客ノード、アーク、旅行時間または距離、二値変数 x\_ij、各顧客を一度訪問する制約、部分巡回除去制約、総旅行時間最小化を定義する。時間依存TSPでは旅行時間が定数 d\_ij ではなく出発時刻 t に依存する d\_ij(t) となり、あるアーク選択が後続の到着時刻と後続アークの旅行時間を変えるため、静的モデルより決定変数・制約・解法が複雑になる。}
\ExamItem{SS24 Aufgabe 1b (4点)}{旅行時間の情報モデルの例と、常に最も詳細なモデルを選ばない理由を問う問題。}{例として、単一平均旅行時間行列、時間帯別行列、区間旅行時間、確率分布モデルがある。詳細なモデルほど現実に近い可能性はあるが、データ量、信頼性、計算時間、解釈可能性、意思決定モデルの複雑性が悪化する。したがって目的に十分な粒度を選ぶ。}
\ExamItem{WS22/23 Aufgabe 4 (8点)}{過去の交通観測から時間帯別平均旅行時間を構築する手順を説明する問題。}{まずFCDなどから位置・時刻・速度を収集する。位置を道路セグメントにマッチングし、セグメントごとの観測速度を得る。外れ値やエラーを除去し、時間帯 w とセグメント l ごとに観測を集約する。平均速度または旅行時間を計算し、観測不足の部分は類似セグメントや時間帯で一般化する。最後に時間帯別の旅行時間行列を作り、TSP/VRPの情報モデルとして使う。}
\ExamItem{SS24 Aufgabe 1b (4点)}{旅行時間情報モデルの選択と複雑性のトレードオフを問う問題。}{平均値モデルは簡単だがラッシュアワーを無視しやすい。時間帯別モデルは実際の交通変動を反映するが、観測数が少ないセルでは信頼性が低くなる。したがってデータ量と意思決定への効果を見ながら集約・一般化する必要がある。}

\NoteSection{確認問題と解答}
\begin{enumerate}
\item \textbf{時間依存旅行時間モデルでは、なぜ一つの旅行時間行列では不十分か。}\\
一つの行列では全時間帯で同じ旅行時間を仮定することになる。実際には道路セグメントごとに混雑時間が異なり、10時と17時では同じ顧客間の最短経路や所要時間が異なる。そのため時間帯別または連続時刻の情報モデルが必要になる。
\item \textbf{時間依存性がヒューリスティックを難しくする理由を説明せよ。}\\
ヒューリスティックの局所的な選択が、後続の出発時刻を変え、後続の旅行時間関数の評価値を変えるからである。静的な追加距離だけでは全体影響を評価できない。
\end{enumerate}

\end{multicols}

\clearpage

\SlideHeader{Lecture 2: Information and Decision Modeling}{030}{セグメントから経路へ}

\SlideImage{30}{../Materials/2026/3DM_02_Information-and-Decision-Modeling_SS26.pdf}

\vspace{1.5mm}

\begin{multicols}{2}

\footnotesize

\raggedcolumns

\NoteSection{日本語訳}

\begin{itemize}
\item 顧客間には複数の潜在的経路がある。
\item 最短経路はDijkstraアルゴリズムで計算される。
\item 時間依存性を考慮する必要がある。つまり時間依存最短経路を探す。
\end{itemize}

\NoteSection{解説}
時間依存旅行時間では、顧客 i から j への旅行時間を定数ではなく出発時刻 t の関数として扱う。記号的には d\_ij(t) または tau\_ij(t) と考える。\par

同じ二地点間でも、出発時刻が異なれば最短経路も旅行時間も変わる。したがって、訪問順序の変更は単にアークの足し算を変えるだけでなく、後続の到着時刻と全体の旅行時間を変える。\par

このため、情報モデルは時間帯別行列や時間依存リンク旅行時間を必要とし、意思決定モデルは時間の次元を含む変数・制約を必要とする。\par

\NoteSection{試験対策ポイント}
旅行時間を d\_ij(t) として捉え、出発時刻が経路・順序・後続旅行時間に与える影響を説明できること。\par

\NoteSection{過去問関連}
\ExamItem{WS22/23 Aufgabe 1 (7点)}{TSP の数理最適化モデルと、時間次元を持つTSPで何が変わるかを問う問題。}{標準TSPでは、顧客ノード、アーク、旅行時間または距離、二値変数 x\_ij、各顧客を一度訪問する制約、部分巡回除去制約、総旅行時間最小化を定義する。時間依存TSPでは旅行時間が定数 d\_ij ではなく出発時刻 t に依存する d\_ij(t) となり、あるアーク選択が後続の到着時刻と後続アークの旅行時間を変えるため、静的モデルより決定変数・制約・解法が複雑になる。}
\ExamItem{SS24 Aufgabe 1b (4点)}{旅行時間の情報モデルの例と、常に最も詳細なモデルを選ばない理由を問う問題。}{例として、単一平均旅行時間行列、時間帯別行列、区間旅行時間、確率分布モデルがある。詳細なモデルほど現実に近い可能性はあるが、データ量、信頼性、計算時間、解釈可能性、意思決定モデルの複雑性が悪化する。したがって目的に十分な粒度を選ぶ。}
\ExamItem{WS22/23 Aufgabe 4 (8点)}{過去の交通観測から時間帯別平均旅行時間を構築する手順を説明する問題。}{まずFCDなどから位置・時刻・速度を収集する。位置を道路セグメントにマッチングし、セグメントごとの観測速度を得る。外れ値やエラーを除去し、時間帯 w とセグメント l ごとに観測を集約する。平均速度または旅行時間を計算し、観測不足の部分は類似セグメントや時間帯で一般化する。最後に時間帯別の旅行時間行列を作り、TSP/VRPの情報モデルとして使う。}
\ExamItem{SS24 Aufgabe 1b (4点)}{旅行時間情報モデルの選択と複雑性のトレードオフを問う問題。}{平均値モデルは簡単だがラッシュアワーを無視しやすい。時間帯別モデルは実際の交通変動を反映するが、観測数が少ないセルでは信頼性が低くなる。したがってデータ量と意思決定への効果を見ながら集約・一般化する必要がある。}

\NoteSection{確認問題と解答}
\begin{enumerate}
\item \textbf{時間依存旅行時間モデルでは、なぜ一つの旅行時間行列では不十分か。}\\
一つの行列では全時間帯で同じ旅行時間を仮定することになる。実際には道路セグメントごとに混雑時間が異なり、10時と17時では同じ顧客間の最短経路や所要時間が異なる。そのため時間帯別または連続時刻の情報モデルが必要になる。
\item \textbf{時間依存性がヒューリスティックを難しくする理由を説明せよ。}\\
ヒューリスティックの局所的な選択が、後続の出発時刻を変え、後続の旅行時間関数の評価値を変えるからである。静的な追加距離だけでは全体影響を評価できない。
\end{enumerate}

\end{multicols}

\clearpage

\SlideHeader{Lecture 2: Information and Decision Modeling}{031}{セグメントから経路へ}

\SlideImage{31}{../Materials/2026/3DM_02_Information-and-Decision-Modeling_SS26.pdf}

\vspace{1.5mm}

\begin{multicols}{2}

\footnotesize

\raggedcolumns

\NoteSection{日本語訳}

\begin{itemize}
\item 異なる時刻に出発すると、異なる道路セグメントを通る可能性がある。
\item 時間依存Dijkstraはこれを考慮する。
\end{itemize}

\NoteSection{解説}
時間依存旅行時間では、顧客 i から j への旅行時間を定数ではなく出発時刻 t の関数として扱う。記号的には d\_ij(t) または tau\_ij(t) と考える。\par

同じ二地点間でも、出発時刻が異なれば最短経路も旅行時間も変わる。したがって、訪問順序の変更は単にアークの足し算を変えるだけでなく、後続の到着時刻と全体の旅行時間を変える。\par

このため、情報モデルは時間帯別行列や時間依存リンク旅行時間を必要とし、意思決定モデルは時間の次元を含む変数・制約を必要とする。\par

\NoteSection{試験対策ポイント}
旅行時間を d\_ij(t) として捉え、出発時刻が経路・順序・後続旅行時間に与える影響を説明できること。\par

\NoteSection{過去問関連}
\ExamItem{WS22/23 Aufgabe 1 (7点)}{TSP の数理最適化モデルと、時間次元を持つTSPで何が変わるかを問う問題。}{標準TSPでは、顧客ノード、アーク、旅行時間または距離、二値変数 x\_ij、各顧客を一度訪問する制約、部分巡回除去制約、総旅行時間最小化を定義する。時間依存TSPでは旅行時間が定数 d\_ij ではなく出発時刻 t に依存する d\_ij(t) となり、あるアーク選択が後続の到着時刻と後続アークの旅行時間を変えるため、静的モデルより決定変数・制約・解法が複雑になる。}
\ExamItem{SS24 Aufgabe 1b (4点)}{旅行時間の情報モデルの例と、常に最も詳細なモデルを選ばない理由を問う問題。}{例として、単一平均旅行時間行列、時間帯別行列、区間旅行時間、確率分布モデルがある。詳細なモデルほど現実に近い可能性はあるが、データ量、信頼性、計算時間、解釈可能性、意思決定モデルの複雑性が悪化する。したがって目的に十分な粒度を選ぶ。}
\ExamItem{WS22/23 Aufgabe 4 (8点)}{過去の交通観測から時間帯別平均旅行時間を構築する手順を説明する問題。}{まずFCDなどから位置・時刻・速度を収集する。位置を道路セグメントにマッチングし、セグメントごとの観測速度を得る。外れ値やエラーを除去し、時間帯 w とセグメント l ごとに観測を集約する。平均速度または旅行時間を計算し、観測不足の部分は類似セグメントや時間帯で一般化する。最後に時間帯別の旅行時間行列を作り、TSP/VRPの情報モデルとして使う。}
\ExamItem{SS24 Aufgabe 1b (4点)}{旅行時間情報モデルの選択と複雑性のトレードオフを問う問題。}{平均値モデルは簡単だがラッシュアワーを無視しやすい。時間帯別モデルは実際の交通変動を反映するが、観測数が少ないセルでは信頼性が低くなる。したがってデータ量と意思決定への効果を見ながら集約・一般化する必要がある。}

\NoteSection{確認問題と解答}
\begin{enumerate}
\item \textbf{時間依存旅行時間モデルでは、なぜ一つの旅行時間行列では不十分か。}\\
一つの行列では全時間帯で同じ旅行時間を仮定することになる。実際には道路セグメントごとに混雑時間が異なり、10時と17時では同じ顧客間の最短経路や所要時間が異なる。そのため時間帯別または連続時刻の情報モデルが必要になる。
\item \textbf{時間依存性がヒューリスティックを難しくする理由を説明せよ。}\\
ヒューリスティックの局所的な選択が、後続の出発時刻を変え、後続の旅行時間関数の評価値を変えるからである。静的な追加距離だけでは全体影響を評価できない。
\end{enumerate}

\end{multicols}

\clearpage

\SlideHeader{Lecture 2: Information and Decision Modeling}{032}{セグメントから経路へ: FIFO}

\SlideImage{32}{../Materials/2026/3DM_02_Information-and-Decision-Modeling_SS26.pdf}

\vspace{1.5mm}

\begin{multicols}{2}

\footnotesize

\raggedcolumns

\NoteSection{日本語訳}

\begin{itemize}
\item 行列の硬い切り替えは奇妙な遷移現象を引き起こす。
\item FIFO: 遅く出発したときに追い越しが起きないようにする必要がある。
\item ブレークポイントを線形結合で滑らかにする。
\end{itemize}

\NoteSection{解説}
FIFO は First-In-First-Out の意味で、同じリンクを先に出発した車両が、後から出発した車両に追い越されて後に到着するようなモデルを避ける条件である。\par

時間帯別行列を階段状に切り替えると、9:59出発では3分、10:00出発では1分のような不連続が起こり、遅く出発した方が早く到着する矛盾が生じる可能性がある。\par

この矛盾を避けるため、ブレークポイント付近を線形結合で平滑化し、出発時刻から到着時刻への関数が単調になるようにする。\par

\NoteSection{試験対策ポイント}
FIFO条件の意味、違反例、平滑化による補正を数値例で説明できること。\par

\NoteSection{過去問関連}
\ExamItem{WS23/24 Aufgabe 4 (6点)}{時間帯別平均旅行時間において FIFO 条件がなぜ必要か、違反時に何が起こるかを説明する問題。}{FIFO は、同じリンクを遅く出発した車両が早く出発した車両を追い越して先に到着してはならないという条件である。時間帯境界で旅行時間を急に切り替えると、9:59出発が長時間、10:00出発が短時間となり、後発が先着する不自然なモデルになり得る。解決には境界付近を線形補間で平滑化し、出発時刻から到着時刻への関数を単調に保つ必要がある。}
\ExamItem{SS25 Exercise 2 (6点)}{FIFO条件を説明し、旅行時間モデルでどのように整合性を保つかを問う問題。}{同じ道路リンク上で出発順と到着順が逆転しないよう、時間依存旅行時間関数のブレークポイントを補正する。時間帯別の平均速度をそのまま階段関数として使うのではなく、境界を滑らかにつなぐ。}

\NoteSection{確認問題と解答}
\begin{enumerate}
\item \textbf{FIFO条件を数値例で説明せよ。}\\
9:59に出発した車両の旅行時間が3分なら到着は10:02である。一方、10:00に出発した車両の旅行時間が1分なら到着は10:01で、後発が先着する。これは同一道路上では不自然でありFIFO違反である。
\item \textbf{FIFO違反を避けるにはどうするか。}\\
時間帯境界で旅行時間を急に切り替えず、境界前後の旅行時間または速度を補間して滑らかにする。これにより出発時刻が遅くなるほど到着時刻も遅くなる単調性を保つ。
\end{enumerate}

\end{multicols}

\clearpage

\SlideHeader{Lecture 2: Information and Decision Modeling}{033}{モデリングの概略}

\SlideImage{33}{../Materials/2026/3DM_02_Information-and-Decision-Modeling_SS26.pdf}

\vspace{1.5mm}

\begin{multicols}{2}

\footnotesize

\raggedcolumns

\NoteSection{日本語訳}

\begin{itemize}
\item 情報モデル
\begin{itemize}
\item 顧客間、デポを含む時間依存旅行時間。
\item 時間依存旅行時間行列、または複数の旅行時間行列。
\end{itemize}
\item 意思決定モデル
\begin{itemize}
\item 目的: 旅行時間を最小化する。
\item 制約: 全顧客をサービスする。
\item 制約: ルーティング制約、部分巡回除去など。
\item 決定変数: 時刻 t に顧客 i と j の間を移動するか。
\end{itemize}
\end{itemize}

\NoteSection{解説}
時間依存旅行時間では、顧客 i から j への旅行時間を定数ではなく出発時刻 t の関数として扱う。記号的には d\_ij(t) または tau\_ij(t) と考える。\par

同じ二地点間でも、出発時刻が異なれば最短経路も旅行時間も変わる。したがって、訪問順序の変更は単にアークの足し算を変えるだけでなく、後続の到着時刻と全体の旅行時間を変える。\par

このため、情報モデルは時間帯別行列や時間依存リンク旅行時間を必要とし、意思決定モデルは時間の次元を含む変数・制約を必要とする。\par

\NoteSection{試験対策ポイント}
旅行時間を d\_ij(t) として捉え、出発時刻が経路・順序・後続旅行時間に与える影響を説明できること。\par

\NoteSection{過去問関連}
\ExamItem{WS22/23 Aufgabe 1 (7点)}{TSP の数理最適化モデルと、時間次元を持つTSPで何が変わるかを問う問題。}{標準TSPでは、顧客ノード、アーク、旅行時間または距離、二値変数 x\_ij、各顧客を一度訪問する制約、部分巡回除去制約、総旅行時間最小化を定義する。時間依存TSPでは旅行時間が定数 d\_ij ではなく出発時刻 t に依存する d\_ij(t) となり、あるアーク選択が後続の到着時刻と後続アークの旅行時間を変えるため、静的モデルより決定変数・制約・解法が複雑になる。}
\ExamItem{SS24 Aufgabe 1b (4点)}{旅行時間の情報モデルの例と、常に最も詳細なモデルを選ばない理由を問う問題。}{例として、単一平均旅行時間行列、時間帯別行列、区間旅行時間、確率分布モデルがある。詳細なモデルほど現実に近い可能性はあるが、データ量、信頼性、計算時間、解釈可能性、意思決定モデルの複雑性が悪化する。したがって目的に十分な粒度を選ぶ。}
\ExamItem{WS22/23 Aufgabe 4 (8点)}{過去の交通観測から時間帯別平均旅行時間を構築する手順を説明する問題。}{まずFCDなどから位置・時刻・速度を収集する。位置を道路セグメントにマッチングし、セグメントごとの観測速度を得る。外れ値やエラーを除去し、時間帯 w とセグメント l ごとに観測を集約する。平均速度または旅行時間を計算し、観測不足の部分は類似セグメントや時間帯で一般化する。最後に時間帯別の旅行時間行列を作り、TSP/VRPの情報モデルとして使う。}
\ExamItem{SS24 Aufgabe 1b (4点)}{旅行時間情報モデルの選択と複雑性のトレードオフを問う問題。}{平均値モデルは簡単だがラッシュアワーを無視しやすい。時間帯別モデルは実際の交通変動を反映するが、観測数が少ないセルでは信頼性が低くなる。したがってデータ量と意思決定への効果を見ながら集約・一般化する必要がある。}

\NoteSection{確認問題と解答}
\begin{enumerate}
\item \textbf{時間依存旅行時間モデルでは、なぜ一つの旅行時間行列では不十分か。}\\
一つの行列では全時間帯で同じ旅行時間を仮定することになる。実際には道路セグメントごとに混雑時間が異なり、10時と17時では同じ顧客間の最短経路や所要時間が異なる。そのため時間帯別または連続時刻の情報モデルが必要になる。
\item \textbf{時間依存性がヒューリスティックを難しくする理由を説明せよ。}\\
ヒューリスティックの局所的な選択が、後続の出発時刻を変え、後続の旅行時間関数の評価値を変えるからである。静的な追加距離だけでは全体影響を評価できない。
\end{enumerate}

\end{multicols}

\clearpage

\SlideHeader{Lecture 2: Information and Decision Modeling}{034}{時間依存決定変数}

\SlideImage{34}{../Materials/2026/3DM_02_Information-and-Decision-Modeling_SS26.pdf}

\vspace{1.5mm}

\begin{multicols}{2}

\footnotesize

\raggedcolumns

\NoteSection{日本語訳}

\begin{itemize}
\item アークごとの旅行時間は時刻により異なる。
\item 例: d\_AB は時刻により 1, 1, 欠落のような値を取る。
\item 時刻を含むネットワーク表現が必要になる。
\end{itemize}

\NoteSection{解説}
時間依存旅行時間では、顧客 i から j への旅行時間を定数ではなく出発時刻 t の関数として扱う。記号的には d\_ij(t) または tau\_ij(t) と考える。\par

同じ二地点間でも、出発時刻が異なれば最短経路も旅行時間も変わる。したがって、訪問順序の変更は単にアークの足し算を変えるだけでなく、後続の到着時刻と全体の旅行時間を変える。\par

このため、情報モデルは時間帯別行列や時間依存リンク旅行時間を必要とし、意思決定モデルは時間の次元を含む変数・制約を必要とする。\par

\NoteSection{試験対策ポイント}
旅行時間を d\_ij(t) として捉え、出発時刻が経路・順序・後続旅行時間に与える影響を説明できること。\par

\NoteSection{過去問関連}
\ExamItem{WS22/23 Aufgabe 1 (7点)}{TSP の数理最適化モデルと、時間次元を持つTSPで何が変わるかを問う問題。}{標準TSPでは、顧客ノード、アーク、旅行時間または距離、二値変数 x\_ij、各顧客を一度訪問する制約、部分巡回除去制約、総旅行時間最小化を定義する。時間依存TSPでは旅行時間が定数 d\_ij ではなく出発時刻 t に依存する d\_ij(t) となり、あるアーク選択が後続の到着時刻と後続アークの旅行時間を変えるため、静的モデルより決定変数・制約・解法が複雑になる。}
\ExamItem{SS24 Aufgabe 1b (4点)}{旅行時間の情報モデルの例と、常に最も詳細なモデルを選ばない理由を問う問題。}{例として、単一平均旅行時間行列、時間帯別行列、区間旅行時間、確率分布モデルがある。詳細なモデルほど現実に近い可能性はあるが、データ量、信頼性、計算時間、解釈可能性、意思決定モデルの複雑性が悪化する。したがって目的に十分な粒度を選ぶ。}
\ExamItem{WS22/23 Aufgabe 4 (8点)}{過去の交通観測から時間帯別平均旅行時間を構築する手順を説明する問題。}{まずFCDなどから位置・時刻・速度を収集する。位置を道路セグメントにマッチングし、セグメントごとの観測速度を得る。外れ値やエラーを除去し、時間帯 w とセグメント l ごとに観測を集約する。平均速度または旅行時間を計算し、観測不足の部分は類似セグメントや時間帯で一般化する。最後に時間帯別の旅行時間行列を作り、TSP/VRPの情報モデルとして使う。}
\ExamItem{SS24 Aufgabe 1b (4点)}{旅行時間情報モデルの選択と複雑性のトレードオフを問う問題。}{平均値モデルは簡単だがラッシュアワーを無視しやすい。時間帯別モデルは実際の交通変動を反映するが、観測数が少ないセルでは信頼性が低くなる。したがってデータ量と意思決定への効果を見ながら集約・一般化する必要がある。}

\NoteSection{確認問題と解答}
\begin{enumerate}
\item \textbf{時間依存旅行時間モデルでは、なぜ一つの旅行時間行列では不十分か。}\\
一つの行列では全時間帯で同じ旅行時間を仮定することになる。実際には道路セグメントごとに混雑時間が異なり、10時と17時では同じ顧客間の最短経路や所要時間が異なる。そのため時間帯別または連続時刻の情報モデルが必要になる。
\item \textbf{時間依存性がヒューリスティックを難しくする理由を説明せよ。}\\
ヒューリスティックの局所的な選択が、後続の出発時刻を変え、後続の旅行時間関数の評価値を変えるからである。静的な追加距離だけでは全体影響を評価できない。
\end{enumerate}

\end{multicols}

\clearpage

\SlideHeader{Lecture 2: Information and Decision Modeling}{035}{時間依存決定変数}

\SlideImage{35}{../Materials/2026/3DM_02_Information-and-Decision-Modeling_SS26.pdf}

\vspace{1.5mm}

\begin{multicols}{2}

\footnotesize

\raggedcolumns

\NoteSection{日本語訳}

\begin{itemize}
\item A, B, C, D と時刻 t1 から t7 による時間展開ネットワーク。
\item ある時刻のノードから将来時刻のノードへ移動することで、時刻依存旅行時間を表現する。
\end{itemize}

\NoteSection{解説}
時間依存旅行時間では、顧客 i から j への旅行時間を定数ではなく出発時刻 t の関数として扱う。記号的には d\_ij(t) または tau\_ij(t) と考える。\par

同じ二地点間でも、出発時刻が異なれば最短経路も旅行時間も変わる。したがって、訪問順序の変更は単にアークの足し算を変えるだけでなく、後続の到着時刻と全体の旅行時間を変える。\par

このため、情報モデルは時間帯別行列や時間依存リンク旅行時間を必要とし、意思決定モデルは時間の次元を含む変数・制約を必要とする。\par

\NoteSection{試験対策ポイント}
旅行時間を d\_ij(t) として捉え、出発時刻が経路・順序・後続旅行時間に与える影響を説明できること。\par

\NoteSection{過去問関連}
\ExamItem{WS22/23 Aufgabe 1 (7点)}{TSP の数理最適化モデルと、時間次元を持つTSPで何が変わるかを問う問題。}{標準TSPでは、顧客ノード、アーク、旅行時間または距離、二値変数 x\_ij、各顧客を一度訪問する制約、部分巡回除去制約、総旅行時間最小化を定義する。時間依存TSPでは旅行時間が定数 d\_ij ではなく出発時刻 t に依存する d\_ij(t) となり、あるアーク選択が後続の到着時刻と後続アークの旅行時間を変えるため、静的モデルより決定変数・制約・解法が複雑になる。}
\ExamItem{SS24 Aufgabe 1b (4点)}{旅行時間の情報モデルの例と、常に最も詳細なモデルを選ばない理由を問う問題。}{例として、単一平均旅行時間行列、時間帯別行列、区間旅行時間、確率分布モデルがある。詳細なモデルほど現実に近い可能性はあるが、データ量、信頼性、計算時間、解釈可能性、意思決定モデルの複雑性が悪化する。したがって目的に十分な粒度を選ぶ。}
\ExamItem{WS22/23 Aufgabe 4 (8点)}{過去の交通観測から時間帯別平均旅行時間を構築する手順を説明する問題。}{まずFCDなどから位置・時刻・速度を収集する。位置を道路セグメントにマッチングし、セグメントごとの観測速度を得る。外れ値やエラーを除去し、時間帯 w とセグメント l ごとに観測を集約する。平均速度または旅行時間を計算し、観測不足の部分は類似セグメントや時間帯で一般化する。最後に時間帯別の旅行時間行列を作り、TSP/VRPの情報モデルとして使う。}
\ExamItem{SS24 Aufgabe 1b (4点)}{旅行時間情報モデルの選択と複雑性のトレードオフを問う問題。}{平均値モデルは簡単だがラッシュアワーを無視しやすい。時間帯別モデルは実際の交通変動を反映するが、観測数が少ないセルでは信頼性が低くなる。したがってデータ量と意思決定への効果を見ながら集約・一般化する必要がある。}

\NoteSection{確認問題と解答}
\begin{enumerate}
\item \textbf{時間依存旅行時間モデルでは、なぜ一つの旅行時間行列では不十分か。}\\
一つの行列では全時間帯で同じ旅行時間を仮定することになる。実際には道路セグメントごとに混雑時間が異なり、10時と17時では同じ顧客間の最短経路や所要時間が異なる。そのため時間帯別または連続時刻の情報モデルが必要になる。
\item \textbf{時間依存性がヒューリスティックを難しくする理由を説明せよ。}\\
ヒューリスティックの局所的な選択が、後続の出発時刻を変え、後続の旅行時間関数の評価値を変えるからである。静的な追加距離だけでは全体影響を評価できない。
\end{enumerate}

\end{multicols}

\clearpage

\SlideHeader{Lecture 2: Information and Decision Modeling}{036}{解を見つける}

\SlideImage{36}{../Materials/2026/3DM_02_Information-and-Decision-Modeling_SS26.pdf}

\vspace{1.5mm}

\begin{multicols}{2}

\footnotesize

\raggedcolumns

\NoteSection{日本語訳}

\begin{itemize}
\item 最適計算は難しく、ヒューリスティックが必要である。
\item 4つのヒューリスティック
\begin{itemize}
\item LANTIME: メタヒューリスティック
\item Nearest Neighbor
\item Cheapest Insertion
\item Savings
\end{itemize}
\item 平均値に基づくTSPも解く。
\end{itemize}

\NoteSection{解説}
時間依存TSP/VRPでは厳密最適化が難しいため、Nearest Neighbor、Cheapest Insertion、Savings、メタヒューリスティックなどの近似解法を使う。\par

Nearest Neighbor は現在地から最も近い顧客を選ぶため時間依存性を比較的扱いやすいが、終盤に遠い顧客が残りやすい。Cheapest Insertion は局所的に最も安い挿入を選ぶが、時間依存性により後続時刻が変わるため、以前の評価が無効になりやすい。Savings は直接配送ツアーを統合するが、統合時の時刻影響の評価が難しい。\par

\NoteSection{試験対策ポイント}
各ヒューリスティックの基本発想と、時間依存性により評価が難しくなる理由を説明できること。\par

\NoteSection{過去問関連}
\ExamItem{WS23/24 Aufgabe 5 (7点)}{Cheapest Insertion が時間依存旅行時間でなぜ難しくなるかを問う問題。}{静的TSPでは顧客挿入の追加費用を局所的に評価しやすい。しかし時間依存の場合、ある顧客を挿入すると後続顧客への到着時刻が変わり、その後のアーク旅行時間も変わる。したがって、局所的に安い挿入が全体では悪化する可能性があり、以前の評価が無効になりやすい。}

\NoteSection{確認問題と解答}
\begin{enumerate}
\item \textbf{Nearest Neighbor, Cheapest Insertion, Savings の違いを説明せよ。}\\
Nearest Neighbor は現在位置から最短の次顧客を選ぶ。Cheapest Insertion は未訪問顧客を既存ツアーの最も安い位置へ挿入する。Savings は各顧客への直接往復ツアーを、節約量が大きい順に統合する。
\item \textbf{時間依存性が Cheapest Insertion を難しくする理由を説明せよ。}\\
顧客を挿入すると、それ以降の訪問時刻がずれる。旅行時間が時刻に依存するため、単に挿入アークのコストだけではなく、後続アークの旅行時間も変化する。したがって局所的な安さが全体最適に結びつきにくい。
\end{enumerate}

\end{multicols}

\clearpage

\SlideHeader{Lecture 2: Information and Decision Modeling}{037}{Nearest Neighbor}

\SlideImage{37}{../Materials/2026/3DM_02_Information-and-Decision-Modeling_SS26.pdf}

\vspace{1.5mm}

\begin{multicols}{2}

\footnotesize

\raggedcolumns

\NoteSection{日本語訳}

\begin{itemize}
\item 現在位置から最も近い顧客を次に選ぶ。
\item 時間依存性を扱うことができる。
\item 通常、終盤に長い移動が残る。
\end{itemize}

\NoteSection{解説}
時間依存TSP/VRPでは厳密最適化が難しいため、Nearest Neighbor、Cheapest Insertion、Savings、メタヒューリスティックなどの近似解法を使う。\par

Nearest Neighbor は現在地から最も近い顧客を選ぶため時間依存性を比較的扱いやすいが、終盤に遠い顧客が残りやすい。Cheapest Insertion は局所的に最も安い挿入を選ぶが、時間依存性により後続時刻が変わるため、以前の評価が無効になりやすい。Savings は直接配送ツアーを統合するが、統合時の時刻影響の評価が難しい。\par

\NoteSection{試験対策ポイント}
各ヒューリスティックの基本発想と、時間依存性により評価が難しくなる理由を説明できること。\par

\NoteSection{過去問関連}
\ExamItem{WS23/24 Aufgabe 5 (7点)}{Cheapest Insertion が時間依存旅行時間でなぜ難しくなるかを問う問題。}{静的TSPでは顧客挿入の追加費用を局所的に評価しやすい。しかし時間依存の場合、ある顧客を挿入すると後続顧客への到着時刻が変わり、その後のアーク旅行時間も変わる。したがって、局所的に安い挿入が全体では悪化する可能性があり、以前の評価が無効になりやすい。}

\NoteSection{確認問題と解答}
\begin{enumerate}
\item \textbf{Nearest Neighbor, Cheapest Insertion, Savings の違いを説明せよ。}\\
Nearest Neighbor は現在位置から最短の次顧客を選ぶ。Cheapest Insertion は未訪問顧客を既存ツアーの最も安い位置へ挿入する。Savings は各顧客への直接往復ツアーを、節約量が大きい順に統合する。
\item \textbf{時間依存性が Cheapest Insertion を難しくする理由を説明せよ。}\\
顧客を挿入すると、それ以降の訪問時刻がずれる。旅行時間が時刻に依存するため、単に挿入アークのコストだけではなく、後続アークの旅行時間も変化する。したがって局所的な安さが全体最適に結びつきにくい。
\end{enumerate}

\end{multicols}

\clearpage

\SlideHeader{Lecture 2: Information and Decision Modeling}{038}{Cheapest Insertion}

\SlideImage{38}{../Materials/2026/3DM_02_Information-and-Decision-Modeling_SS26.pdf}

\vspace{1.5mm}

\begin{multicols}{2}

\footnotesize

\raggedcolumns

\NoteSection{日本語訳}

\begin{itemize}
\item 空のツアーから始める。
\item 段階的に最も安い顧客をツアーに挿入する。
\item 顧客がずらされ、旅行時間の変化により以前の決定が有効でなくなる。
\end{itemize}

\NoteSection{解説}
時間依存TSP/VRPでは厳密最適化が難しいため、Nearest Neighbor、Cheapest Insertion、Savings、メタヒューリスティックなどの近似解法を使う。\par

Nearest Neighbor は現在地から最も近い顧客を選ぶため時間依存性を比較的扱いやすいが、終盤に遠い顧客が残りやすい。Cheapest Insertion は局所的に最も安い挿入を選ぶが、時間依存性により後続時刻が変わるため、以前の評価が無効になりやすい。Savings は直接配送ツアーを統合するが、統合時の時刻影響の評価が難しい。\par

\NoteSection{試験対策ポイント}
各ヒューリスティックの基本発想と、時間依存性により評価が難しくなる理由を説明できること。\par

\NoteSection{過去問関連}
\ExamItem{WS23/24 Aufgabe 5 (7点)}{Cheapest Insertion が時間依存旅行時間でなぜ難しくなるかを問う問題。}{静的TSPでは顧客挿入の追加費用を局所的に評価しやすい。しかし時間依存の場合、ある顧客を挿入すると後続顧客への到着時刻が変わり、その後のアーク旅行時間も変わる。したがって、局所的に安い挿入が全体では悪化する可能性があり、以前の評価が無効になりやすい。}

\NoteSection{確認問題と解答}
\begin{enumerate}
\item \textbf{Nearest Neighbor, Cheapest Insertion, Savings の違いを説明せよ。}\\
Nearest Neighbor は現在位置から最短の次顧客を選ぶ。Cheapest Insertion は未訪問顧客を既存ツアーの最も安い位置へ挿入する。Savings は各顧客への直接往復ツアーを、節約量が大きい順に統合する。
\item \textbf{時間依存性が Cheapest Insertion を難しくする理由を説明せよ。}\\
顧客を挿入すると、それ以降の訪問時刻がずれる。旅行時間が時刻に依存するため、単に挿入アークのコストだけではなく、後続アークの旅行時間も変化する。したがって局所的な安さが全体最適に結びつきにくい。
\end{enumerate}

\end{multicols}

\clearpage

\SlideHeader{Lecture 2: Information and Decision Modeling}{039}{Savings}

\SlideImage{39}{../Materials/2026/3DM_02_Information-and-Decision-Modeling_SS26.pdf}

\vspace{1.5mm}

\begin{multicols}{2}

\footnotesize

\raggedcolumns

\NoteSection{日本語訳}

\begin{itemize}
\item 直接往復ツアーから始める。
\item 最小費用でツアーを統合する。
\item 時間依存性の評価は難しい。
\item 少なくとも後のステップでは暗黙に考慮される。
\end{itemize}

\NoteSection{解説}
時間依存TSP/VRPでは厳密最適化が難しいため、Nearest Neighbor、Cheapest Insertion、Savings、メタヒューリスティックなどの近似解法を使う。\par

Nearest Neighbor は現在地から最も近い顧客を選ぶため時間依存性を比較的扱いやすいが、終盤に遠い顧客が残りやすい。Cheapest Insertion は局所的に最も安い挿入を選ぶが、時間依存性により後続時刻が変わるため、以前の評価が無効になりやすい。Savings は直接配送ツアーを統合するが、統合時の時刻影響の評価が難しい。\par

\NoteSection{試験対策ポイント}
各ヒューリスティックの基本発想と、時間依存性により評価が難しくなる理由を説明できること。\par

\NoteSection{過去問関連}
\ExamItem{WS23/24 Aufgabe 5 (7点)}{Cheapest Insertion が時間依存旅行時間でなぜ難しくなるかを問う問題。}{静的TSPでは顧客挿入の追加費用を局所的に評価しやすい。しかし時間依存の場合、ある顧客を挿入すると後続顧客への到着時刻が変わり、その後のアーク旅行時間も変わる。したがって、局所的に安い挿入が全体では悪化する可能性があり、以前の評価が無効になりやすい。}

\NoteSection{確認問題と解答}
\begin{enumerate}
\item \textbf{Nearest Neighbor, Cheapest Insertion, Savings の違いを説明せよ。}\\
Nearest Neighbor は現在位置から最短の次顧客を選ぶ。Cheapest Insertion は未訪問顧客を既存ツアーの最も安い位置へ挿入する。Savings は各顧客への直接往復ツアーを、節約量が大きい順に統合する。
\item \textbf{時間依存性が Cheapest Insertion を難しくする理由を説明せよ。}\\
顧客を挿入すると、それ以降の訪問時刻がずれる。旅行時間が時刻に依存するため、単に挿入アークのコストだけではなく、後続アークの旅行時間も変化する。したがって局所的な安さが全体最適に結びつきにくい。
\end{enumerate}

\end{multicols}

\clearpage

\SlideHeader{Lecture 2: Information and Decision Modeling}{040}{結果}

\SlideImage{40}{../Materials/2026/3DM_02_Information-and-Decision-Modeling_SS26.pdf}

\vspace{1.5mm}

\begin{multicols}{2}

\footnotesize

\raggedcolumns

\NoteSection{日本語訳}

\begin{itemize}
\item 出発時刻が異なる。
\item 旅行時間は時刻に依存する。
\item 静的モデルは過大評価も過小評価も行う。
\end{itemize}

\NoteSection{解説}
時間依存旅行時間では、顧客 i から j への旅行時間を定数ではなく出発時刻 t の関数として扱う。記号的には d\_ij(t) または tau\_ij(t) と考える。\par

同じ二地点間でも、出発時刻が異なれば最短経路も旅行時間も変わる。したがって、訪問順序の変更は単にアークの足し算を変えるだけでなく、後続の到着時刻と全体の旅行時間を変える。\par

このため、情報モデルは時間帯別行列や時間依存リンク旅行時間を必要とし、意思決定モデルは時間の次元を含む変数・制約を必要とする。\par

\NoteSection{試験対策ポイント}
旅行時間を d\_ij(t) として捉え、出発時刻が経路・順序・後続旅行時間に与える影響を説明できること。\par

\NoteSection{過去問関連}
\ExamItem{WS22/23 Aufgabe 1 (7点)}{TSP の数理最適化モデルと、時間次元を持つTSPで何が変わるかを問う問題。}{標準TSPでは、顧客ノード、アーク、旅行時間または距離、二値変数 x\_ij、各顧客を一度訪問する制約、部分巡回除去制約、総旅行時間最小化を定義する。時間依存TSPでは旅行時間が定数 d\_ij ではなく出発時刻 t に依存する d\_ij(t) となり、あるアーク選択が後続の到着時刻と後続アークの旅行時間を変えるため、静的モデルより決定変数・制約・解法が複雑になる。}
\ExamItem{SS24 Aufgabe 1b (4点)}{旅行時間の情報モデルの例と、常に最も詳細なモデルを選ばない理由を問う問題。}{例として、単一平均旅行時間行列、時間帯別行列、区間旅行時間、確率分布モデルがある。詳細なモデルほど現実に近い可能性はあるが、データ量、信頼性、計算時間、解釈可能性、意思決定モデルの複雑性が悪化する。したがって目的に十分な粒度を選ぶ。}
\ExamItem{WS22/23 Aufgabe 4 (8点)}{過去の交通観測から時間帯別平均旅行時間を構築する手順を説明する問題。}{まずFCDなどから位置・時刻・速度を収集する。位置を道路セグメントにマッチングし、セグメントごとの観測速度を得る。外れ値やエラーを除去し、時間帯 w とセグメント l ごとに観測を集約する。平均速度または旅行時間を計算し、観測不足の部分は類似セグメントや時間帯で一般化する。最後に時間帯別の旅行時間行列を作り、TSP/VRPの情報モデルとして使う。}
\ExamItem{SS24 Aufgabe 1b (4点)}{旅行時間情報モデルの選択と複雑性のトレードオフを問う問題。}{平均値モデルは簡単だがラッシュアワーを無視しやすい。時間帯別モデルは実際の交通変動を反映するが、観測数が少ないセルでは信頼性が低くなる。したがってデータ量と意思決定への効果を見ながら集約・一般化する必要がある。}

\NoteSection{確認問題と解答}
\begin{enumerate}
\item \textbf{時間依存旅行時間モデルでは、なぜ一つの旅行時間行列では不十分か。}\\
一つの行列では全時間帯で同じ旅行時間を仮定することになる。実際には道路セグメントごとに混雑時間が異なり、10時と17時では同じ顧客間の最短経路や所要時間が異なる。そのため時間帯別または連続時刻の情報モデルが必要になる。
\item \textbf{時間依存性がヒューリスティックを難しくする理由を説明せよ。}\\
ヒューリスティックの局所的な選択が、後続の出発時刻を変え、後続の旅行時間関数の評価値を変えるからである。静的な追加距離だけでは全体影響を評価できない。
\end{enumerate}

\end{multicols}

\clearpage

\SlideHeader{Lecture 2: Information and Decision Modeling}{041}{実装}

\SlideImage{41}{../Materials/2026/3DM_02_Information-and-Decision-Modeling_SS26.pdf}

\vspace{1.5mm}

\begin{multicols}{2}

\footnotesize

\raggedcolumns

\NoteSection{日本語訳}

\begin{itemize}
\item 月曜正午に出発
\begin{itemize}
\item 最初は市中心部を避ける。
\item 先に郊外をサービスする。
\end{itemize}
\item 月曜夕方に出発
\begin{itemize}
\item 先に市中心部をサービスする。
\item 最初に郊外道路を使う。
\item 異なる経路になる。
\end{itemize}
\end{itemize}

\NoteSection{解説}
時間依存旅行時間では、顧客 i から j への旅行時間を定数ではなく出発時刻 t の関数として扱う。記号的には d\_ij(t) または tau\_ij(t) と考える。\par

同じ二地点間でも、出発時刻が異なれば最短経路も旅行時間も変わる。したがって、訪問順序の変更は単にアークの足し算を変えるだけでなく、後続の到着時刻と全体の旅行時間を変える。\par

このため、情報モデルは時間帯別行列や時間依存リンク旅行時間を必要とし、意思決定モデルは時間の次元を含む変数・制約を必要とする。\par

\NoteSection{試験対策ポイント}
旅行時間を d\_ij(t) として捉え、出発時刻が経路・順序・後続旅行時間に与える影響を説明できること。\par

\NoteSection{過去問関連}
\ExamItem{WS22/23 Aufgabe 1 (7点)}{TSP の数理最適化モデルと、時間次元を持つTSPで何が変わるかを問う問題。}{標準TSPでは、顧客ノード、アーク、旅行時間または距離、二値変数 x\_ij、各顧客を一度訪問する制約、部分巡回除去制約、総旅行時間最小化を定義する。時間依存TSPでは旅行時間が定数 d\_ij ではなく出発時刻 t に依存する d\_ij(t) となり、あるアーク選択が後続の到着時刻と後続アークの旅行時間を変えるため、静的モデルより決定変数・制約・解法が複雑になる。}
\ExamItem{SS24 Aufgabe 1b (4点)}{旅行時間の情報モデルの例と、常に最も詳細なモデルを選ばない理由を問う問題。}{例として、単一平均旅行時間行列、時間帯別行列、区間旅行時間、確率分布モデルがある。詳細なモデルほど現実に近い可能性はあるが、データ量、信頼性、計算時間、解釈可能性、意思決定モデルの複雑性が悪化する。したがって目的に十分な粒度を選ぶ。}
\ExamItem{WS22/23 Aufgabe 4 (8点)}{過去の交通観測から時間帯別平均旅行時間を構築する手順を説明する問題。}{まずFCDなどから位置・時刻・速度を収集する。位置を道路セグメントにマッチングし、セグメントごとの観測速度を得る。外れ値やエラーを除去し、時間帯 w とセグメント l ごとに観測を集約する。平均速度または旅行時間を計算し、観測不足の部分は類似セグメントや時間帯で一般化する。最後に時間帯別の旅行時間行列を作り、TSP/VRPの情報モデルとして使う。}
\ExamItem{SS24 Aufgabe 1b (4点)}{旅行時間情報モデルの選択と複雑性のトレードオフを問う問題。}{平均値モデルは簡単だがラッシュアワーを無視しやすい。時間帯別モデルは実際の交通変動を反映するが、観測数が少ないセルでは信頼性が低くなる。したがってデータ量と意思決定への効果を見ながら集約・一般化する必要がある。}

\NoteSection{確認問題と解答}
\begin{enumerate}
\item \textbf{時間依存旅行時間モデルでは、なぜ一つの旅行時間行列では不十分か。}\\
一つの行列では全時間帯で同じ旅行時間を仮定することになる。実際には道路セグメントごとに混雑時間が異なり、10時と17時では同じ顧客間の最短経路や所要時間が異なる。そのため時間帯別または連続時刻の情報モデルが必要になる。
\item \textbf{時間依存性がヒューリスティックを難しくする理由を説明せよ。}\\
ヒューリスティックの局所的な選択が、後続の出発時刻を変え、後続の旅行時間関数の評価値を変えるからである。静的な追加距離だけでは全体影響を評価できない。
\end{enumerate}

\end{multicols}

\clearpage

\SlideHeader{Lecture 2: Information and Decision Modeling}{042}{要するに}

\SlideImage{42}{../Materials/2026/3DM_02_Information-and-Decision-Modeling_SS26.pdf}

\vspace{1.5mm}

\begin{multicols}{2}

\footnotesize

\raggedcolumns

\NoteSection{日本語訳}

\begin{itemize}
\item 問題とデータを真剣に扱うには多くの作業が必要である。
\item 情報モデル: データ収集、分析、構造化、集約、一般化。
\item 意思決定モデル: より多くの決定変数、複雑な意思決定モデル、適応した解法。
\item 得られるアプローチはより詳細なモデリングを行うが、決定論的挙動を仮定している。
\end{itemize}

\NoteSection{解説}
Real-World は、配送、交通、顧客、車両、作業員など、実際に行動が行われる対象である。ここでは全ての詳細をモデルに入れることはできないため、意思決定に必要な情報だけを選ぶ。\par

Aggregation は、観測データをモデルで使える次元へ変換する処理である。例えばGPS位置列を道路セグメント別の平均速度へ変換すること、住所を顧客ノードへ変換することが該当する。\par

Information Model は、現実を圧縮した情報表現である。旅行時間行列、時間帯別速度パターン、顧客集合、需要分布などが含まれる。\par

Decision Model は、可能な決定、目的関数、制約、評価方法を定義する。TSP/VRPであれば、訪問順序、利用アーク、各顧客を一度訪問する制約、総旅行時間最小化が含まれる。\par

Implementation は、モデルの解を現実の行動へ戻す処理である。例えば『1-3-2』というノード順序を、実際の住所と道路ナビゲーションへ変換する。\par

\NoteSection{試験対策ポイント}
Real-World, Aggregation, Information Model, Decision Model, Implementation の役割とデータ・解の流れを、配送例に対応付けて説明できること。\par

\NoteSection{過去問関連}
\ExamItem{WS23/24 Aufgabe 2 (5点)}{Real-World Problem と Decision Model の関係、および Information Model を介した流れを説明する問題。}{現実世界から直接最適化モデルを作るのではなく、まず観測データを集約して意思決定に必要な情報表現に落とす。Information Model は現実を圧縮した表現で、Decision Model はその表現を使って決定空間、制約、目的関数を定義する。解は Implementation により現実の行動へ変換され、結果観測から仮説・モデルが更新される。}
\ExamItem{SS24 Aufgabe 1a (8点)}{Information Model, Decision Model, Aggregation, Disaggregation/Implementation の機能を説明する問題。}{Aggregation は現実データをモデルの次元へ変換する。Information Model は意思決定に必要な状態・属性・パラメータを整理する。Decision Model は可能な決定、目的、制約、評価方法を定義する。Disaggregation/Implementation はモデル解を現実の作業指示やルートへ戻す。重要なのは、これらが一方向で終わらず、実装後の観測で更新されるループである点である。}
\ExamItem{WS22/23 Aufgabe 3 (6点)}{Business Analytics プロセスの中で OR 的なモデル化・解法がどこにあるかを説明する問題。}{情報モデルはデータ側、意思決定モデルはOR側の橋渡しである。ORは、現実の問題構造を抽象化し、目的関数・制約・決定変数に変換し、解を評価可能にする。}
\ExamItem{WS22/23 Aufgabe 2 (6点)}{Business Analytics を構成する分野と、それらの交差領域を説明する問題。}{Operations Research は意思決定モデル、制約、目的関数、最適化を担う。Statistics はデータの不確実性、分布、推定、予測を扱う。Computer Science/Information Systems はデータ管理、アルゴリズム、実装、情報システムを担う。Business Analytics はこれらを組み合わせ、データから意思決定可能な知見と実行可能な解を作る分野である。}
\ExamItem{WS22/23 Aufgabe 3 (6点)}{Business Analytics の適用プロセスを挙げ、Operations Research に属する工程を示す問題。}{現実問題を観察し、データを取得・集約して情報モデルを作り、そこから意思決定モデルを定義し、解法で解を求め、実装して現実に戻す。ORに特に属するのは、目的関数・制約・決定変数を持つ意思決定モデルの定式化、解空間の探索、最適化・ヒューリスティックによる解の計算である。}
\ExamItem{WS23/24 Aufgabe 1 (6点)}{Business Analytics の三つの構成分野と交差領域を説明する問題。}{三分野を単に列挙するだけでなく、データ分析だけでは意思決定にならず、最適化だけでも現実データを反映できない点を述べる。交差領域では、データからモデルを作り、モデルから現実に実装可能な意思決定を作ることが中心になる。}
\ExamItem{WS23/24 Aufgabe 3 (10点)}{配送車両の業務的ツアープランニングに Business Analytics の構成要素を対応付ける問題。}{Real-world は配送現場、顧客、道路、車両である。Aggregation は住所・道路・走行時間観測をモデル変数へ変換すること。Information model は顧客集合、距離/旅行時間行列、時間帯情報である。Decision model は訪問順序・経路・目的関数・制約を持つVRP/TSPである。Implementation は得られたツアーをドライバーのルート案内として実行することである。}

\NoteSection{確認問題と解答}
\begin{enumerate}
\item \textbf{Information Model と Decision Model の違いを説明せよ。}\\
Information Model は、意思決定に必要な現実情報の表現である。例として顧客、道路、旅行時間行列、時間帯別速度がある。Decision Model は、その情報を用いて決定を選ぶための目的関数、制約、決定変数の体系である。情報モデルが『何を知っているか』を表し、意思決定モデルが『その情報で何をどう決めるか』を表す。
\item \textbf{Aggregation と Implementation の役割を説明せよ。}\\
Aggregation は現実データを抽象化してモデル入力に変換する処理であり、Implementation はモデル解を実際の行動へ具体化する処理である。前者は現実からモデルへの橋、後者はモデルから現実への橋である。
\item \textbf{なぜこのループは一回で終わらないのか。}\\
実装後に現実の結果が観測されるためである。予測より配送が遅れた、交通パターンが違った、制約が現実に合わなかった場合、仮説、情報モデル、意思決定モデルを更新する必要がある。
\end{enumerate}

\end{multicols}

\clearpage

\SlideHeader{Lecture 2: Information and Decision Modeling}{043}{次回予告}

\SlideImage{43}{../Materials/2026/3DM_02_Information-and-Decision-Modeling_SS26.pdf}

\vspace{1.5mm}

\begin{multicols}{2}

\footnotesize

\raggedcolumns

\NoteSection{日本語訳}

\begin{itemize}
\item 不確実性と、それが情報モデル・意思決定モデルに与える影響。
\end{itemize}

\end{multicols}